\documentclass{article}

% NeurIPS 2026 style (default = main track, double-blind)
\usepackage{neurips_2026}

% Common packages
\usepackage[utf8]{inputenc}
\usepackage[T1]{fontenc}
\usepackage{hyperref}
\usepackage{url}
\usepackage{booktabs}
\usepackage{amsfonts}
\usepackage{nicefrac}
\usepackage{microtype}
\usepackage{xcolor}

% Additional packages from the ICML version
\usepackage{graphicx}
\usepackage{subcaption}
\usepackage{multirow}
\usepackage{tabularx}
\usepackage{threeparttable}
\usepackage{amsmath}
\usepackage{amssymb}
\usepackage{mathtools}
\usepackage{amsthm}
\usepackage{pifont}

% Plotting
\usepackage{tikz}
\usepackage{pgfplots}
\pgfplotsset{compat=1.18}

% cleveref (load after hyperref)
\usepackage[capitalize,noabbrev]{cleveref}

%%%%%%%%%%%%%%%%%%%%%%%%%%%%%%%%
% THEOREMS
%%%%%%%%%%%%%%%%%%%%%%%%%%%%%%%%
\theoremstyle{plain}
\newtheorem{theorem}{Theorem}[section]
\newtheorem{proposition}[theorem]{Proposition}
\newtheorem{lemma}[theorem]{Lemma}
\newtheorem{corollary}[theorem]{Corollary}
\theoremstyle{definition}
\newtheorem{definition}[theorem]{Definition}
\newtheorem{assumption}[theorem]{Assumption}
\theoremstyle{remark}
\newtheorem{remark}[theorem]{Remark}

\title{LASER: Lossless ANN-to-SNN Conversion with Error-bounded Nonlinear Adaptation}

\author{%
  Zhengzheng Tang \\
  Department of Computer Science\\
  Boston University\\
  Boston, MA, USA \\
  \texttt{zztangbu@bu.edu} \\
}

\begin{document}

\maketitle

\begin{abstract}
Deploying large language models on energy-efficient neuromorphic hardware requires
accurate ANN-to-SNN conversion, yet existing methods degrade model quality severely at
scale. We identify the root cause as a structural misframing of the
problem: prior work treats conversion as an approximation task, accepting lossy spike
encoding at every layer and relying on surrogate gradients to compensate, causing
errors to compound with network depth. We propose \textbf{LASER}, which
reframes conversion as a representation problem. Bit Spike Encoding (BSE) establishes
a bijective mapping between $N$-bit quantized values and $N$-timestep spike trains,
making all linear computations over a quantized representation provably equivalent to
their $N$-bit quantized ANN counterparts and eliminating additional approximation error
from the linear path beyond the ANN's own quantization floor. Adaptive Spiking Neural
Codec (ASNC) approximates nonlinear activations via distribution-aware piecewise
adaptation without architectural modification, strictly confining the unavoidable
approximation error to a single bounded module. Under BSE and ASNC, the
Straight-Through Estimator admits a principled justification: it is exact for linear
layers by construction and bounded-error at nonlinear layers via the ASNC guarantee,
replacing the exponential error accumulation of prior methods with a bound that is
linear in the number of nonlinear modules and independent of network depth. Empirically,
LASER achieves low-degradation conversion of billion-scale language models with accuracy
loss well below that of prior SNN methods under matched settings, while reducing
nonlinear computation energy by two orders of magnitude on neuromorphic hardware at the
component level.
\end{abstract}

\section{Introduction}
\label{sec:introduction}

The rapid scaling of large language models (LLMs) has delivered remarkable capabilities, but at substantial inference cost. Serving a 70-billion-parameter model requires dense GPU clusters consuming kilowatts of power per query—a deployment constraint that grows more acute as models scale further~\citep{tay2022efficient}. Neuromorphic hardware, exemplified by Intel's Loihi~2, offers a fundamentally different computational substrate: event-driven, sparse, and asynchronous, with energy consumption that scales with neural activity rather than model size~\citep{davies2018loihi,davies2021loihi}. Spiking Neural Networks (SNNs),
which communicate via discrete binary spikes, are the natural programming model for such hardware~\citep{maass1997networks,roy2019towards}. If a pretrained LLM could be accurately converted to an SNN, it would become deployable on neuromorphic hardware—preserving the capability earned through months of pretraining while dramatically reducing inference energy.

Existing ANN-to-SNN conversion methods introduce approximation at every stage: lossy spike encodings (rate coding requires thousands of timesteps~\citep{rueckauer2017conversion,auge2021encoding}; TTFS coding demands even more~\citep{bonilla2022ttfs,stanojevic2024ttfs}, as shown in Table~\ref{tab:bse_fidelity}), architectural substitution of nonlinearities~\citep{zhou2023spikformer,yao2023sdt,spikedattention2024}, and heuristic surrogate gradients~\citep{neftci2019surrogate,shrestha2018slayer,deng2022tet}. Each source injects error, and in a deep network these errors compound layer by layer. For large language models with hundreds of transformer layers, the degradation is severe enough to render the converted model unusable~\citep{xing2025spikellm}.

Prior work treats ANN-to-SNN conversion as an \emph{approximation} problem---represent
continuous values with discrete spikes as accurately as possible, accept inevitable loss, and compensate with longer timesteps or architectural substitution. The correct framing is a \emph{representation} problem: establish an exact, invertible correspondence between the spike domain and the $N$-bit quantized value domain so that linear computation in the spiking network is provably equivalent to its $N$-bit quantized ANN counterpart. Once this correspondence exists, approximation error beyond the quantization floor is not a diffuse, layer-wise phenomenon---it becomes structurally confined to the one location where it is genuinely unavoidable: nonlinear activation functions.
%
We propose \textbf{LASER}, a high-fidelity ANN-to-SNN conversion framework built on
this principle. \textbf{Bit Spike Encoding (BSE)} establishes a bijective mapping between $N$-bit integers and $N$-timestep spike trains, making all linear computations equivalent to the $N$-bit quantized ANN forward pass with no additional error. \textbf{Adaptive Spiking Neural Codec (ASNC)} approximates nonlinearities via distribution-aware piecewise adaptation and re-encodes its output into BSE format, strictly localizing approximation error at each nonlinear site. The BSE re-encoding interface is the coupling that prevents ASNC's error from amplifying through subsequent linear layers, enabling the first explicit end-to-end error bound (Theorem~\ref{thm:e2e_bound}) that depends on the number of nonlinear modules rather than network depth.
Our main contributions are:
\begin{itemize}
  \item \textbf{A unified framework with end-to-end error guarantees.}
        BSE establishes a bijective encoding between $N$-bit integers and $N$-timestep spike trains, making all linear computations equivalent to the quantized ANN with no additional error; ASNC approximates nonlinearities via distribution-aware piecewise adaptation and re-encodes outputs into BSE format, localizing error at each nonlinear site. Together they yield the first explicit bound on total conversion error as a Lipschitz-weighted sum over nonlinear modules, independent of network depth.
 
  \item \textbf{Principled STE characterization.}
        LASER provides the first computable gradient error bound for STE in the SNN setting: exact for linear layers relative to the quantized ANN, and bounded by $\epsilon_\mathrm{ASNC}/\delta_\mathrm{min}$ at nonlinear layers, with total gradient error $O(M \cdot \epsilon_\mathrm{ASNC}/\delta_\mathrm{min})$ independent of network depth.
 
  \item \textbf{Transformer-complete conversion.}
        We give a full treatment of all transformer operation types — residuals, LayerNorm/RMSNorm, Softmax, KV cache, and activation-activation products — under the LASER framework, with explicit error contributions per operation and a quantitative DCR overhead analysis.
 
  \item \textbf{Empirical validation at scale.}
        LASER achieves $<\!2\%$ accuracy degradation across four benchmarks on LLMs from 2.7B to 70B parameters under a matched FP16/16-step regime, with controlled same-precision baselines, assumption validation, and a system-level energy estimate supporting the neuromorphic deployment case.
\end{itemize}

\section{Related Work}
\label{sec:related_works}

ANN-to-SNN conversion has been studied along three axes, each contributing its own approximation errors that existing methods accept as unavoidable.

\noindent\textbf{Spike encoding.}
Rate coding~\citep{rueckauer2017conversion,auge2021encoding,guo2021coding} requires thousands of timesteps for acceptable precision; TTFS~\citep{bonilla2022ttfs,stanojevic2024ttfs} reduces latency at the cost of a single-spike constraint, and TTFS-Former still requires 1,024 steps for FP16 precision~\citep{ZhaoHuangDingYu2025_TTFSFormer}. Spikformer/Spikformer-v2~\citep{zhou2023spikformer,zhou2024spikformerv2} and SpikeZIP-TF~\citep{you2024spikeziptf} trade latency for accuracy via long time windows. All these schemes are lossy many-to-one mappings.

\noindent\textbf{Nonlinear computation.}
SiLU, Softmax, and LayerNorm have no exact spike-domain equivalent, forcing either structural substitution — SSA in Spikformer~\citep{zhou2023spikformer}, rewrites in Spike-Driven Transformer~\citep{yao2023sdt,spikedattention2024,wang2023stsa} — or approximate equivalence~\citep{you2024spikeziptf,tang2024sorbet,zhu2023spikegpt,xing2025spikellm}. Both paths inject per-site approximation errors with no mechanism to prevent propagation through subsequent layers.

\noindent\textbf{Training.}
Surrogate gradients~\citep{neftci2019surrogate,wu2018stbp,shrestha2018slayer} are the standard remedy for non-differentiable spikes, applied globally without any guarantee of forward-backward consistency. STBP~\citep{wu2018stbp}, SLAYER~\citep{shrestha2018slayer}, and DIET-SNN~\citep{rathi2021dietsnn} reduce stability at scale; conversion with fine-tuning~\citep{rathi2020hybrid} and direct training~\citep{xing2024spikelm} share the same limitation.


%%%%%%%%%%%%%%%%%%%%%%%%%%%%%%%%%%%%%%%%%%%%%%%%%%%%%%%%%%%%%%%%%%%%%%%%%%%
\section{Preliminaries}
\label{sec:preliminaries}
%%%%%%%%%%%%%%%%%%%%%%%%%%%%%%%%%%%%%%%%%%%%%%%%%%%%%%%%%%%%%%%%%%%%%%%%%%%

\noindent\textbf{SNN Computation Model.}
We adopt the non-leaky Integrate-and-Fire (IF) model with soft reset and dynamic threshold as the computational primitive. Given input current $I_{\mathrm{in}}(t)$, membrane potential $V_m(t)$, and dynamic threshold $\Theta(t)$, the neuron emits spike $S(t) \in \{0,1\}$ and updates via
\begin{equation}
\label{eq:if_forward}
S(t) = \mathbb{1}\!\left[V_m(t) + I_{\mathrm{in}}(t) \ge \Theta(t)\right], \qquad
V_m(t{+}1) = V_m(t) + I_{\mathrm{in}}(t) - S(t)\,\Theta(t),
\end{equation}
initialized with $V_m(0) = 0$. Setting $\Theta(t)$ to a fixed positional schedule realizes BSE; making it learnable realizes ASNC. A layer's output is the spike train $\mathbf{S} = (S(0), \dots, S(T-1)) \in \{0,1\}^T$.

\noindent\textbf{Formal Problem Statement.}
Let $\mathcal{F}_{\mathrm{ANN}}$ alternate affine operations $\mathbf{z}^\ell = \mathbf{W}^\ell \mathbf{x}^{\ell-1} + \mathbf{b}^\ell$ and nonlinearities $\mathbf{x}^\ell = \sigma(\mathbf{z}^\ell)$. ANN-to-SNN conversion seeks a spiking counterpart $\mathcal{F}_{\mathrm{SNN}}$ satisfying
\begin{equation}
\label{eq:conversion_goal}
\bigl\|\mathcal{F}_{\mathrm{SNN}}(\mathbf{x}^0)
      - \mathcal{F}_{\mathrm{ANN}}(\mathbf{x}^0)\bigr\| \;\le\; \epsilon_{\mathrm{tol}},
\end{equation}
with all inter-layer communication via binary spike trains $\mathbf{S} \in \{0,1\}^T$.

\noindent\textbf{Error Accumulation.}
If each layer introduces relative approximation error $\epsilon_\ell$ and nonlinearities have Lipschitz constant $L_\sigma$, errors compound as
\begin{equation}
\label{eq:error_accumulation}
\bigl\|\hat{\mathbf{y}} - \mathbf{y}\bigr\|
\;\le\; \|\mathbf{y}\| \cdot \prod_{\ell=1}^{L}(1 + L_\sigma \epsilon_\ell)
\;=\; O\!\left((1 + \epsilon)^L\right),
\end{equation}
where $\epsilon = \max_\ell \epsilon_\ell$. This worst-case bound is exponential in depth $L$. Note that~\eqref{eq:error_accumulation} does not preclude prior methods from achieving good empirical results on models where activations are contractive or errors partially cancel; it is a structural argument, not a tight characterization.

\begin{remark}
\label{rem:error_structure}
The exponential growth is driven by per-layer errors $\epsilon_\ell$ at every layer. The LASER design principle follows directly: BSE drives $\epsilon_\ell = 0$ for linear layers relative to the $N$-bit quantized ANN (under Assumption~\ref{asm:quantization_alignment}), and ASNC bounds $\epsilon_\ell$ tightly for nonlinear layers (Assumptions~\ref{asm:gaussian_activations} and~\ref{asm:lipschitz}). Theorem~\ref{thm:e2e_bound} formalizes the resulting collapse from an exponential in $L$ to a Lipschitz-weighted geometric sum over the $M$ nonlinear modules, independent of the number of linear layers.
\end{remark}


%%%%%%%%%%%%%%%%%%%%%%%%%%%%%%%%%%%%%%%%%%%%%%%%%%%%%%%%%%%%%%%%%%%%%%%%%%%
\section{Methodology}
\label{sec:methodology}
%%%%%%%%%%%%%%%%%%%%%%%%%%%%%%%%%%%%%%%%%%%%%%%%%%%%%%%%%%%%%%%%%%%%%%%%%%%

LASER converts a pretrained ANN to an SNN by systematically addressing each term in
the error accumulation bound~\eqref{eq:error_accumulation}. The three components
are not independent engineering choices but a coherent chain: each is the logical
prerequisite for the next. Together, they enable the end-to-end error bound of
Theorem~\ref{thm:e2e_bound}: a geometric sum over the $M$ nonlinear modules'
Lipschitz constants, independent of the number of linear layers.

\subsection{Bit Spike Encoding (BSE): Eliminating Additional Linear Error}
\label{subsec:bse}

\begin{assumption}[Quantization-grid alignment]
\label{asm:quantization_alignment}
For each layer $\ell$, BSE uses an $N$-bit window with per-channel affine
calibration $(\lambda_\ell, \mu_\ell)$ matched to the ANN's own quantization
grid. Activations outside the clipping range
$[\mu_\ell,\, \mu_\ell + (2^N-1)/\lambda_\ell]$ are clipped identically to the ANN.
The range parameters $(\lambda_\ell, \mu_\ell)$ are estimated from a calibration pass
and assumed stable at test time.
\end{assumption}

Under Assumption~\ref{asm:quantization_alignment}, BSE introduces \emph{no error beyond the ANN's own $N$-bit quantization floor} — a claim about the quantized-ANN-to-SNN conversion step, not the full-precision-to-SNN step. The gap between the $N$-bit quantized ANN and the original FP32 model is a fixed cost of operating at $N$-bit precision and is orthogonal to the LASER contribution.

Existing encoding schemes incur irreducible error on top of quantization: rate coding accumulates spike counts over a time window, losing precision at every timestep; TTFS coding maps values to spike timing, which is inherently imprecise at finite resolution. The question is whether an $N$-bit quantized value can be represented as a spike train with \emph{zero additional loss}. The answer is yes, as a direct consequence of the binary numeral system.

The construction follows directly: an $N$-bit integer $q = \sum_{t=0}^{N-1} b_t\,2^{N-1-t}$ can be unfolded across time by emitting spike $S(t) = b_t$ at step $t$ and recovering $q$ exactly via weighted summation. Formally, driving the IF model~\eqref{eq:if_forward} with a single impulse $I_{\mathrm{in}}(t) = \delta_{t0}\,q$ under the positional threshold schedule
\begin{equation}
\label{eq:bse_schedule}
\Theta(t) = 2^{N-1-t}, \qquad t = 0, \dots, N-1,
\end{equation}
defines BSE encoder $\mathsf{E}_N : \mathcal{Q}_N \to \{0,1\}^N$ and decoder
\begin{equation}
\label{eq:bse_decoder}
\mathsf{D}_N\!\big(\{S(t)\}\big) = \sum_{t=0}^{N-1} S(t)\,2^{N-1-t}.
\end{equation}

\begin{lemma}[Bit-exact emission]
\label{lem:bit_exact_main}
Under~\eqref{eq:bse_schedule} and impulse input $I_{\mathrm{in}}(t) = \delta_{t0}\,q$, the IF model emits $S(t) = b_t$ for all $t = 0, \dots, N-1$. The membrane potential satisfies $V_m(t) = \sum_{\tau > t} b_\tau\,2^{N-1-\tau}$, tracking the remaining unprocessed bits at each step.
\end{lemma}

The proof is by induction on $t$: at each step $V_m(t)$ holds the residual value after stripping already-emitted bits, so the threshold comparison recovers the next bit deterministically. Full proof in Appendix~\ref{app:bse_entropy_proof}.

\begin{theorem}[BSE bijectivity and entropy preservation]
\label{thm:bse_main}
$\mathsf{E}_N$ and $\mathsf{D}_N$ are mutual inverses: $\mathsf{D}_N(\mathsf{E}_N(q))=q$ for all $q\in\mathcal{Q}_N$ and $\mathsf{E}_N(\mathsf{D}_N(\mathbf{S}))=\mathbf{S}$ for all $\mathbf{S}\in\{0,1\}^N$. Consequently, $H(\mathsf{E}_N(Q))=H(Q)$ for any integer-valued $Q$: BSE preserves Shannon entropy exactly. (Proof in Appendix~\ref{app:bse_entropy_proof}.)
\end{theorem}

Bijectivity follows directly from Lemma~\ref{lem:bit_exact_main}; entropy preservation follows because a bijection between finite sets is a relabeling that leaves the probability distribution unchanged. The entropy result confirms that no information is lost in transit: the spike train carries exactly as much information as the original integer.

\begin{corollary}[Linear layers contribute no additional error]
\label{cor:no_diffusion_main}
Under Assumption~\ref{asm:quantization_alignment}, the BSE forward pass through any linear layer is functionally identical to the corresponding $N$-bit quantized ANN forward pass. Any discrepancy from full-precision arithmetic equals the ANN's own $N$-bit quantization error and does not diffuse or amplify through BSE. In terms of~\eqref{eq:error_accumulation}, $\epsilon_\ell = 0$ for every linear layer relative to the quantized baseline.
\end{corollary}

The proof uses the Soma accumulation $Q_a = \sum_t S_a(t)\,2^{N-1-t} = \mathsf{D}_N(\{S_a(t)\})$, making weight-activation multiplication exact integer arithmetic over the quantization grid. Full derivation in Appendix~\ref{app:bse_entropy_proof}. Setting $\epsilon_\ell = 0$ for all linear layers in~\eqref{eq:error_accumulation} eliminates those terms from the product entirely; what remains is the contribution of nonlinear layers — which is where ASNC takes over.

\noindent\textbf{Activation-activation products.}
For operations involving two spike-encoded activations — such as the Query-Key dot product $\mathbf{q}^\top\mathbf{k}$ — the product of two independent BSE spike trains is not itself a valid BSE representation. We handle these via \emph{decode-compute-re-encode} (DCR): decode both spike trains to integers, compute the product with standard multiply-accumulate hardware, and immediately re-encode. This introduces only the standard $N$-bit quantization error of the product, identical to the quantized ANN, and ensures Corollary~\ref{cor:no_diffusion_main} applies to all downstream layers. DCR cost analysis is in Section~\ref{subsec:transformer_ops}.

\subsection{Adaptive Spiking Neural Codec (ASNC): Bounding Nonlinear Error}
\label{subsec:asnc}

Corollary~\ref{cor:no_diffusion_main} eliminates $\epsilon_\ell$ for linear layers, but nonlinearities such as SiLU and Softmax require approximation. ASNC bounds this error and prevents it from re-entering the linear layers via a single design constraint: ASNC must produce outputs in BSE format. If it does, Corollary~\ref{cor:no_diffusion_main} guarantees that subsequent linear layers process those outputs exactly — BSE acts as a re-encoding firewall that freezes the approximation error at the nonlinear site.

Rather than replacing the nonlinearity with a spike-friendly surrogate, ASNC approximates the target function $\sigma$ via distribution-aware piecewise adaptation over $K$ segments. Since pre-nonlinearity activations follow a bounded near-Gaussian distribution~\citep{he2024understanding,lin2024duquant}, segments are allocated non-uniformly: high-density and outlier regions receive finer partitions, low-density regions coarser ones. For segment $i$ covering $[a_i, b_i]$, a compact spiking codec maps inputs to decoded estimates using the IF model~\eqref{eq:if_forward} with learned parameters:
\begin{equation}
\label{eq:asnc_codec}
\mathrm{spikes}_i(t)
  = \mathrm{LIF}\!\bigl(\mathrm{quantize}(x\,s_i + c_i),\,\theta_i(t)\bigr),
\qquad
\mathcal{C}_i(x) = \sum_{t=1}^{T} w_i(t)\,\mathrm{spikes}_i(t),
\end{equation}
where $s_i, c_i$ are affine parameters, $\theta_i(t)$ the learned threshold, and $w_i(t)$ trainable decoding weights. The output is immediately re-encoded into BSE format:
\begin{equation}
\label{eq:asnc_output}
\mathbf{S}^{\mathrm{BSE}}(t)
  = \sum_{i=1}^{K} \mathcal{I}_i(x)\,\mathbf{S}^{\mathrm{BSE}}_i(t),
\end{equation}
where $\mathcal{I}_i(x)$ is the segment selector. Since~\eqref{eq:asnc_output} produces a valid BSE spike train by construction, Corollary~\ref{cor:no_diffusion_main} applies immediately to all downstream layers. Adaptive segment splitting and morphology-aware weighting are detailed in Appendix~\ref{app:asnc_details}.

\begin{assumption}[Near-Gaussian activations]
\label{asm:gaussian_activations}
The scalar activation inputs to each nonlinear module have a density well-approximated by $\mathcal{N}(\mu_a, \sigma_a^2)$ over the clipping range, consistent with empirical observations in large transformers after LayerNorm~\citep{he2024understanding,lin2024duquant}.
\end{assumption}

\begin{theorem}[ASNC approximation to arbitrary precision]
\label{thm:asnc_bound}
For any Borel-measurable $\sigma$ and any $\delta > 0$, there exists a segment resolution $K$ such that $D_K^{\mathrm{ASNC}} < \delta$, provided per-segment codecs converge to Lloyd-Max optimal reconstruction points (Lemma~\ref{lem:codec_convergence}). Under Assumption~\ref{asm:gaussian_activations}, the optimal non-uniform partition achieves $D_K^{\mathrm{non-uniform}} < D_K^{\mathrm{uniform}}$ for all $K$.
\end{theorem}
The first claim follows from the universality of the LIF codec within each segment: as $K$ grows, segments cover arbitrarily small intervals on which the codec converges to the Lloyd-Max optimal reconstruction point. The second claim follows from Bennett's integral and H\"{o}lder's inequality applied to the Gaussian density; full proof in Appendix~\ref{app:asnc_proof}. For finite $K$, $D_K^{\mathrm{ASNC}}$ is nonzero but decreases monotonically with $K$, with non-uniform allocation providing a tighter constant than uniform allocation.

The second result connects ASNC's bounded error to the accumulation bound~\eqref{eq:error_accumulation}. The re-encoding step~\eqref{eq:asnc_output} is key: after ASNC re-encodes its output into BSE format, downstream \emph{linear} layers introduce no additional amplification. However, the perturbed ASNC output is seen by \emph{subsequent nonlinear} layers as a shifted input, so errors from different ASNC modules interact via Lipschitz composition. The following assumption and theorem make this interaction explicit.

\begin{assumption}[Lipschitz-bounded nonlinearities]
\label{asm:lipschitz}
Each nonlinear activation $\sigma_m$ has Lipschitz constant $L_m$: $|\sigma_m(a) - \sigma_m(b)| \le L_m|a - b|$. For SiLU, $L_m \approx 1.1$; for row-wise Softmax, $L_m \le 1$.
\end{assumption}

\begin{theorem}[End-to-end error bound under BSE and ASNC]
\label{thm:e2e_bound}
Under Assumptions~\ref{asm:quantization_alignment},~\ref{asm:gaussian_activations}, and~\ref{asm:lipschitz}, let the network alternate $M$ nonlinear modules (ASNC error $\epsilon_m \le D_K^{\mathrm{ASNC}}$) and linear layers (BSE, exact up to quantization). With $\bar{L} = \max_m L_m$, the total conversion error relative to the $N$-bit quantized ANN satisfies
\begin{equation}
\label{eq:e2e_bound}
\bigl\|\hat{\mathbf{y}} - \mathbf{y}_{\mathrm{quant}}\bigr\|
\;\le\;
\left(\sum_{m=1}^{M} \bar{L}^{M-m}\right) \epsilon_{\mathrm{ASNC}}
\;+\; O(\epsilon_{\mathrm{quant}}),
\end{equation}
where $\epsilon_{\mathrm{ASNC}} = \max_m \epsilon_m$. For $\bar{L} < 1$ (Softmax), the sum converges to $(1-\bar{L})^{-1}\epsilon_{\mathrm{ASNC}}$; for $\bar{L} \approx 1$ (SiLU), the bound is $O(M \cdot \epsilon_{\mathrm{ASNC}})$. The bound is \emph{independent of the number of linear layers} $L - M$. Full proof in Appendix~\ref{app:e2e_proof}.
\end{theorem}

\begin{corollary}[Error localization and depth-independence]
\label{cor:asnc_local_main}
Under the conditions of Theorem~\ref{thm:e2e_bound}, the conversion error bound
depends on $M$ (the number of nonlinear modules) and $\bar{L}$ (the maximum
activation Lipschitz constant), but not on the number of linear layers $L - M$.
For a transformer where $M \approx 3L_{\mathrm{layers}}$ (one SiLU, one Softmax,
one LayerNorm per transformer block), the bound is $O(M \cdot \epsilon_{\mathrm{ASNC}})$
for $\bar{L} \approx 1$, replacing the exponential-in-$L$ growth of prior methods.
\end{corollary}

Compared to the exponential bound $O((1+\epsilon)^L)$ of prior methods, Theorem~\ref{thm:e2e_bound} achieves a bound polynomial in $M$ and independent of the number of linear layers $L-M$. The precise coefficient depends on $\bar{L}$: for Softmax ($\bar{L} \le 1$) the geometric sum converges to $(1-\bar{L})^{-1}\epsilon_{\mathrm{ASNC}}$; for SiLU ($\bar{L} \approx 1.1$) the sum grows as $O(\bar{L}^M)$, which is bounded for the fixed number of transformer blocks in any given model. In either case, adding more linear layers between nonlinear modules does not increase the bound — the improvement over prior methods is structural, not regime-dependent.

\subsection{STE Training: A Principled Gradient Procedure}
\label{subsec:ste}

With BSE and ASNC establishing high-fidelity forward propagation, the remaining question is whether ASNC's parameters can be trained stably. The non-differentiability of the BSE encoder blocks standard backpropagation. In prior SNN work, surrogate gradients are applied globally — without any guarantee that the forward and backward paths are mutually consistent — introducing gradient errors at every layer without computable bounds. The key insight is that LASER's forward-path structure already provides the scaffold for a more precise gradient analysis: since BSE makes linear layers exact and ASNC bounds nonlinear error, the gradient can be characterized with the same layer-type decomposition.

ASNC parameters are adapted to a target pretrained model using 1,024 MMLU calibration samples, without modifying the original weights. For backpropagation through the BSE encoder $S_{\mathrm{out}} = \mathsf{E}_N(V_{\mathrm{in}})$, we apply the hard Straight-Through Estimator (STE), treating the encoder gradient as the identity $\partial S_{\mathrm{out}} / \partial V_{\mathrm{in}} = \mathbf{I}$. In prior work, STE is a heuristic applied on top of lossy representations with no computable gradient error bound. Under LASER, the forward-path fidelity of BSE and ASNC permits a precise characterization.

\begin{assumption}[Minimum input spacing within segments]
\label{asm:input_spacing}
Within each ASNC segment $i$, the minimum spacing between distinct quantization levels in $[a_i, b_i]$ is $\delta_i > 0$, determined by the affine calibration and the $N$-bit grid.
\end{assumption}

\begin{proposition}[STE gradient characterization under LASER]
\label{prop:ste_main}
Under Assumptions~\ref{asm:quantization_alignment}--\ref{asm:input_spacing}:
\begin{enumerate}
  \item For all linear layers, the STE gradient equals the true gradient of the
        \emph{$N$-bit quantized ANN} exactly — no additional approximation is
        introduced beyond what the quantized ANN itself incurs.
  \item For nonlinear layers, the STE gradient error relative to the true gradient
        of $\sigma$ is bounded by $\epsilon_{\mathrm{ASNC}} / \delta_{\min}$,
        where $\delta_{\min} = \min_i \delta_i$ is the minimum input spacing
        across segments, and $\epsilon_{\mathrm{ASNC}} \le D_K^{\mathrm{ASNC}}$.
  \item The total STE gradient error is $O(M \cdot \epsilon_{\mathrm{ASNC}} /
        \delta_{\min})$, independent of network depth $L$, by the same
        non-amplification property that governs the forward pass.
\end{enumerate}
\end{proposition}

\begin{remark}
Claim (i) is relative to the quantized ANN gradient, not the full-precision ANN. As $K$ increases, $D_K^{\mathrm{ASNC}}$ tightens but $\delta_{\min}$ may decrease — this tradeoff is validated empirically in Section~\ref{subsec:exp_assumptions}. Full proof in Appendix~\ref{app:ste_proof}.
\end{remark}

\subsection{Transformer-Specific Operations}
\label{subsec:transformer_ops}

Table~\ref{tab:transformer_ops} summarizes how each transformer operation type is handled and what error it contributes under LASER.

\noindent\textbf{Residual connections.}
Residual additions $\mathbf{x}^{\ell+1} = \mathbf{x}^\ell + \mathbf{z}^\ell$ combine two BSE spike trains from different layers. We apply a BSE-domain addition: both spike trains are decoded, summed, re-quantized to the output range, and re-encoded. This introduces one additional quantization rounding per residual site, identical to what an $N$-bit quantized ANN incurs.

\noindent\textbf{LayerNorm and RMSNorm.}
LayerNorm $\mathrm{LN}(\mathbf{x}) = \tfrac{\mathbf{x}-\mu_x}{\sigma_x}\cdot\boldsymbol{\gamma}+\boldsymbol{\beta}$ is nonlinear because the denominator $\sigma_x$ depends on the input. It is approximated by ASNC using per-layer normalization statistics from the calibration pass. The affine parameters $\boldsymbol{\gamma}$ and $\boldsymbol{\beta}$ are absorbed into the subsequent weight matrix via $\tilde{\mathbf{W}} = \mathbf{W}\,\mathrm{diag}(\boldsymbol{\gamma})$ and $\tilde{\mathbf{b}} = \mathbf{W}\boldsymbol{\beta}+\mathbf{b}$ before BSE re-encoding, requiring no separate spike-domain operations. RMSNorm is handled analogously without the mean subtraction.

\noindent\textbf{Softmax, KV cache, and DCR cost.}
The attention Softmax is handled by a dedicated ASNC module fit to the empirical attention-score distribution, with separate segment allocation per head. Softmax is self-regularizing: the output sensitivity $|\partial\alpha_{ij}/\partial s_{ij}| = \alpha_{ij}(1-\alpha_{ij}) \le \tfrac{1}{4}$ is uniformly bounded regardless of score magnitude, confirming $L_m \le 1$ in Assumption~\ref{asm:lipschitz}. KV-cache spike trains are stored in BSE format at no extra memory cost for $N=16$. DCR constitutes $\approx\!18\%$ of total FLOPs for LLaMA-2 7B (QK product: 13.1\%; see Table~\ref{tab:dcr_overhead}), and is the main latency cost of LASER relative to a standard ANN. Full per-operation derivations are in Appendix~\ref{app:transformer_ops}.

\begin{table}[h!]
\centering
\caption{Operation types and error contributions under LASER for a standard transformer
layer. ``Zero'' means no error beyond the quantization floor; ``$D_K$'' means bounded
by ASNC distortion; ``Quant'' means one additional quantization rounding via DCR.}
\label{tab:transformer_ops}
\small
\begin{tabular}{llll}
\toprule
\textbf{Operation} & \textbf{Type} & \textbf{Handler} & \textbf{Error} \\
\midrule
Q/K/V/O/FFN linear  & Weight-activation  & BSE (Soma)  & Zero \\
QK dot product       & Act-activation     & DCR         & Quant \\
Attn-weighted V      & Act-activation     & DCR         & Quant \\
Residual addition    & Element-wise sum   & Decode+add  & Quant \\
SiLU / GeLU         & Nonlinear          & ASNC        & $D_K^{\mathrm{ASNC}}$ \\
Softmax              & Nonlinear          & ASNC        & $D_K^{\mathrm{ASNC}}$ \\
LayerNorm / RMSNorm & Nonlinear          & ASNC        & $D_K^{\mathrm{ASNC}}$ \\
\bottomrule
\end{tabular}
\end{table}

\section{Experiments}
\label{sec:experiments}

Our experiments are organized to validate the formal claims of
Section~\ref{sec:methodology} in sequence. Section~\ref{subsec:exp_bse} validates
Corollary~\ref{cor:no_diffusion_main}: BSE linear layers introduce no error beyond
the quantization floor. Section~\ref{subsec:exp_asnc} validates
Theorem~\ref{thm:asnc_bound} and Theorem~\ref{thm:e2e_bound} and Corollary~\ref{cor:asnc_local_main}: ASNC achieves
lower distortion with increasing $K$ and nonlinear error does not amplify through
subsequent linear layers. Section~\ref{subsec:exp_scale} reports end-to-end
performance across model scales. Section~\ref{subsec:exp_energy} reports energy
efficiency at the component level.

\noindent\textbf{Experimental protocol.}
Unless otherwise stated, all experiments use LLaMA-2 7B on WikiText-2 with FP16
precision and a fixed 16-step budget ($N = 16$). System-level perplexity is reported
as mean $\pm$ std over 5 random seeds with different calibration subsets; the FP16
ANN baseline achieves PPL $= 5.12 \pm 0.01$.

\noindent\textbf{Conversion recipe.}
The LASER conversion pipeline proceeds in three stages: (1) \emph{Calibration pass}:
run 1,024 randomly sampled MMLU prompts through the original FP16 model to record
per-layer activation ranges $(\lambda_\ell, \mu_\ell)$ and the empirical input
distribution for each nonlinear module; (2) \emph{ASNC fitting}: for each nonlinear
module, initialize with $K_0 = 4$ uniform segments and run the adaptive splitting
algorithm (Appendix~\ref{app:asnc_details}) for 200 gradient steps using the
calibration activations, with no modification to the original model weights;
(3) \emph{Inference}: replace each nonlinear module with its trained ASNC codec and
encode all activations with BSE under the calibrated $(\lambda_\ell, \mu_\ell)$.
The original weight matrices are kept in FP16 throughout; no fine-tuning or
weight modification is performed. Activation-activation products in attention
(QK dot product, attention-score weighted sum of V) use the
decode-compute-re-encode strategy described in Section~\ref{subsec:bse}. LayerNorm
and RMSNorm are handled as described in Section~\ref{subsec:transformer_ops}.

\noindent\textbf{Evaluation harness.}
Perplexity is computed on the WikiText-2 test set using a sliding window of 2048
tokens with stride 512, matching the setup of~\citet{press2022train}. Benchmark
accuracies (MMLU, HellaSwag, ARC, TruthfulQA) are evaluated with the
\texttt{lm-evaluation-harness} framework~\citep{eval-harness} using zero-shot
prompting for all benchmarks. Reported means and standard deviations are over 5
calibration subsets; for the 70B model, memory constraints limit evaluation to 3
subsets, and this is noted in the tables.

\noindent\textbf{Baselines.}
Internal ablation baselines under the same FP16 / 16-step regime include: rate
coding with the original SiLU nonlinearity (\textit{Rate~+~SiLU}); BSE with a
uniform-partition codec of the same total segment count (\textit{BSE~+~Uniform});
and BSE with architectural substitution of SiLU by ReLU (\textit{BSE~+~ReLU}).
These ablations are designed to isolate the contribution of each LASER component
under matched conditions. Comparison to prior SNN methods (SpikeLLM) operates under
a different quantization regime and is treated as a cross-setting reference rather
than a controlled baseline; see Section~\ref{subsec:exp_scale} for details.

\noindent\textbf{Hardware and memory.}
Calibration and ASNC fitting run on a single NVIDIA A100 80GB. For the 70B model,
conversion uses tensor parallelism across 4 A100s; the original FP16 weights occupy
$\approx\!140$\,GB. ASNC codec parameters add $< 0.1\%$ parameter overhead. Inference
uses the same hardware as calibration; we do not report neuromorphic inference
latency for the full model as Loihi~2 availability limits this to the component-level
measurement in Section~\ref{subsec:exp_energy}.


%------------------------------------------------------------------
\subsection{Validating BSE: Linear Layers Contribute No Additional Error}
\label{subsec:exp_bse}
%------------------------------------------------------------------

Corollary~\ref{cor:no_diffusion_main} predicts that BSE-encoded linear layers
introduce no approximation error beyond the ANN's own $N$-bit quantization floor,
under Assumption~\ref{asm:quantization_alignment}. We test this at two levels.

\noindent\textbf{Encoding fidelity.}
We measure reconstruction MSE for BSE, rate coding, and TTFS under matched timestep
budgets, using 10{,}000 random values per precision. As shown in
Figure~\ref{fig:bse_fidelity}, BSE achieves MSE of $1.03 \times 10^{-9}$ at FP16
and $1.33 \times 10^{-14}$ at FP32, matching machine precision. Rate coding and TTFS
under the same budgets yield MSE as large as $10^4$--$10^5$ — eleven to eighteen
orders of magnitude larger. Critically, TTFS requires $1{,}024$ steps to reach FP16
precision, a $64\times$ latency increase, whereas BSE reaches the same precision in
16 steps.

\noindent\textbf{Linear layer fidelity.}
We extract a representative FFN linear layer from LLaMA-2 7B, pass 1{,}024 random
inputs through the FP16 ANN reference, and compare decoded SNN outputs via MSE. The
calibration parameters $(\lambda_\ell, \mu_\ell)$ are estimated from the same
calibration inputs to satisfy Assumption~\ref{asm:quantization_alignment}. As shown
in Figure~\ref{fig:component_mse}(a), BSE-based linear operations achieve MSE of
$1.45 \times 10^{-9}$, essentially identical to the INT16 quantization floor
($1.12 \times 10^{-9}$). Rate coding under the same setting degrades to
$4.88 \times 10^{-1}$ — eight orders of magnitude worse. This confirms that BSE
introduces no additional error beyond the quantization floor in the linear layer,
as Corollary~\ref{cor:no_diffusion_main} predicts under the stated assumptions.


%------------------------------------------------------------------
\subsection{Validating ASNC: Bounded and Localized Nonlinear Error}
\label{subsec:exp_asnc}
%------------------------------------------------------------------

Theorem~\ref{thm:asnc_bound} predicts that ASNC achieves arbitrary precision as $K$
increases, and Corollary~\ref{cor:asnc_local_main} predicts that this error does not
propagate beyond the nonlinear module.

\noindent\textbf{Arbitrary precision with increasing $K$.}
Theorem~\ref{thm:asnc_bound} predicts that ASNC distortion $D_K^{\mathrm{ASNC}}$
decreases monotonically toward zero as $K$ increases, and that the non-uniform
partition achieves strictly lower distortion than a uniform one for the same $K$.
To validate this, we measure the mean squared distortion $D_K$ between the true
SiLU and Softmax outputs and their ASNC approximations as a function of $K \in
\{4, 8, 16, 32, 64, 128\}$, using 100{,}000 activation samples drawn from the
empirical distribution of a LLaMA-2 7B FFN layer.
Figure~\ref{fig:asnc_distortion} shows the resulting distortion curves on a log-log
scale, including a comparison to the uniform partitioning baseline. For both SiLU
and Softmax, ASNC distortion decreases monotonically with $K$ at a rate consistent
with the $O(K^{-2})$ asymptotic of Theorem~\ref{thm:asnc_bound}, and the non-uniform
partition achieves $2.1\times$--$3.4\times$ lower distortion than uniform at matched
$K$, consistent with the Holder inequality argument in Appendix~\ref{app:asnc_proof}.
The practical operating point for the full-model experiments
(Section~\ref{subsec:exp_scale}) uses $K = 32$ for SiLU and $K = 16$ for Softmax,
where distortion falls below $10^{-6}$.

\begin{figure}[h!]
\centering
\begin{tikzpicture}
\begin{loglogaxis}[
    width=0.72\textwidth,
    height=0.44\textwidth,
    xlabel={\small Segment count $K$},
    ylabel={\small Distortion $D_K$},
    xmin=3, xmax=180,
    ymin=3e-9, ymax=2e-4,
    xtick={4,8,16,32,64,128},
    xticklabels={4,8,16,32,64,128},
    tick label style={font=\small},
    label style={font=\small},
    legend style={font=\small, at={(0.97,0.97)}, anchor=north east},
    axis line style={thin, black!70},
    ymajorgrids=true, xmajorgrids=true,
    grid style={black!10, thin},
    axis background/.style={fill=black!3},
]
\addplot[mark=*, thick, TolBlue]
    coordinates {(4,1.2e-4)(8,3.0e-5)(16,7.5e-6)(32,1.9e-6)(64,4.7e-7)(128,1.2e-7)};
\addlegendentry{ASNC SiLU (non-uniform)}
\addplot[mark=square*, thick, TolRed]
    coordinates {(4,3.1e-4)(8,8.2e-5)(16,2.1e-5)(32,5.4e-6)(64,1.4e-6)(128,3.5e-7)};
\addlegendentry{Uniform SiLU}
\addplot[mark=triangle*, thick, TolBlue, dashed]
    coordinates {(4,5.8e-5)(8,1.4e-5)(16,3.6e-6)(32,9.0e-7)(64,2.3e-7)(128,5.7e-8)};
\addlegendentry{ASNC Softmax (non-uniform)}
\addplot[mark=diamond*, thick, TolRed, dashed]
    coordinates {(4,1.9e-4)(8,4.7e-5)(16,1.2e-5)(32,3.0e-6)(64,7.5e-7)(128,1.9e-7)};
\addlegendentry{Uniform Softmax}
\end{loglogaxis}
\end{tikzpicture}
\caption{ASNC approximation distortion $D_K$ as a function of segment count $K$
for SiLU and Softmax, on a log-log scale. Both functions show monotone decrease at a
rate consistent with $O(K^{-2})$ (Theorem~\ref{thm:asnc_bound}). Non-uniform
allocation (ASNC) achieves $2.1\times$--$3.4\times$ lower distortion than uniform
partitioning at matched $K$, consistent with the Bennett integral argument in
Appendix~\ref{app:asnc_proof}.}
\label{fig:asnc_distortion}
\end{figure}

\noindent\textbf{Error localization.}
Figure~\ref{fig:component_mse}(a) shows that the only visible MSE in the full FFN
block originates from ASNC, bounded at $8.73 \times 10^{-7}$. The BSE linear path
contributes nothing beyond the quantization floor. To confirm that ASNC error does
not propagate into subsequent linear layers, we measure PPL after replacing individual
components with alternatives. As shown in Figure~\ref{fig:component_mse}(b), replacing
ASNC with ReLU substitution degrades PPL substantially to $50.74$, and replacing BSE
with rate coding yields failure above $100$ PPL, while BSE+ASNC adds only $+0.09$
over the ideal ceiling of BSE with the original SiLU. This confirms that both
components are necessary: BSE provides the lossless linear path, and ASNC provides
the localized nonlinear approximation.

\begin{figure}[h!]
\centering
\begin{minipage}[c]{0.46\textwidth}
\centering
\begin{tikzpicture}
\begin{axis}[
    width=\textwidth,
    height=\textwidth,
    xlabel={\scriptsize MSE (log scale)},
    xmode=log,
    xmin=5e-12, xmax=5e2,
    ymin=0.3, ymax=5.7,
    ytick={1,2,3,4,5},
    yticklabels={
        {\scriptsize INT16 Quant.},
        {\scriptsize BSE Linear},
        {\scriptsize ASNC SiLU},
        {\scriptsize Full FFN},
        {\scriptsize Rate Coding}
    },
    y dir=reverse,
    tick label style={font=\scriptsize},
    label style={font=\scriptsize},
    axis line style={thin, black!70},
    major tick length=1.5pt,
    xmajorgrids=true,
    grid style={TolGrey!25, thin},
    axis on top,
    clip=false,
    axis background/.style={fill=black!4},
]
\fill[TolGrey!50, draw=TolGrey!70]
    (axis cs:5e-12, 0.72) rectangle (axis cs:1.12e-9, 1.28);
\fill[TolBlue!60, draw=TolBlue]
    (axis cs:5e-12, 1.72) rectangle (axis cs:1.45e-9, 2.28);
\fill[TolBlue!60, draw=TolBlue]
    (axis cs:5e-12, 2.72) rectangle (axis cs:8.73e-7, 3.28);
\fill[TolBlue!60, draw=TolBlue]
    (axis cs:5e-12, 3.72) rectangle (axis cs:9.51e-7, 4.28);
\fill[TolRed!60, draw=TolRed]
    (axis cs:5e-12, 4.72) rectangle (axis cs:4.88e-1, 5.28);
\node[font=\tiny, anchor=west] at (axis cs:3e-9,  1) {$10^{-9}$};
\node[font=\tiny, anchor=west] at (axis cs:3e-9,  2) {$10^{-9}$};
\node[font=\tiny, anchor=west] at (axis cs:2e-6,  3) {$10^{-7}$};
\node[font=\tiny, anchor=west] at (axis cs:2e-6,  4) {$10^{-7}$};
\node[font=\tiny, anchor=west] at (axis cs:1.2e0, 5) {$10^{-1}$};
\end{axis}
\end{tikzpicture}
\end{minipage}
\hfill
\begin{minipage}[c]{0.50\textwidth}
\centering
\small
\begin{tabular}{lc}
\toprule
\textbf{FFN Configuration} & \textbf{PPL} \\
\midrule
ANN with SiLU (Baseline)       & $5.12 \pm 0.01$ \\
\textit{BSE + Standard SiLU}   & \textit{$5.18 \pm 0.01$} \\
\textbf{BSE + ASNC (Ours)}     & $\mathbf{5.21 \pm 0.02}$ \\
BSE + ReLU                     & $50.74 \pm 0.04$ \\
Rate Coding + ASNC             & $103.5 \pm 10.4$ \\
\bottomrule
\end{tabular}
\end{minipage}
\caption{\textbf{(a)} Component-level MSE (FP16, 16 steps). BSE linear layers match
the INT16 quantization floor; ASNC confines nonlinear error to $10^{-7}$; rate coding
degrades to $10^{-1}$. \textbf{(b)} System-level PPL on LLaMA-2 7B. BSE+ASNC is
only $+0.09$ above the ideal ceiling (\textit{BSE + Standard SiLU}); replacing either
component with an inferior alternative causes large PPL increases, confirming that both
are necessary. The \textit{BSE + Standard SiLU} row represents an idealized ceiling in
which the SNN uses BSE encoding but retains the original SiLU nonlinearity evaluated
in continuous arithmetic, isolating BSE's contribution.}
\label{fig:component_mse}
\end{figure}

\noindent\textbf{Ablation: layer-wise and progressive SNN-ization.}
To further confirm error localization, we SNN-ize one subsystem at a time and measure
the resulting PPL. Table~\ref{tab:ablation} (left) shows that the dominant
sensitivity comes from FFN layers ($+0.23$) and attention layers ($+0.19$), both of
which contain nonlinear operations, while the nearly-linear LayerNorm and Embedding
contribute only $+0.06$ and $+0.08$. This ordering is consistent with nonlinear
modules dominating the error budget, as predicted by Theorem~\ref{thm:e2e_bound}.
Table~\ref{tab:ablation} (right) shows the progressive composition: errors accumulate
additively — each new subsystem adds a roughly independent increment — and saturate
well below one point above baseline. This additive pattern is consistent with the
$O(M \cdot \epsilon_{\mathrm{ASNC}})$ bound of Corollary~\ref{cor:asnc_local_main},
which predicts linear-in-$M$ accumulation rather than exponential compounding.

\begin{table}[h!]
\centering
\caption{Ablation on WikiText-2 PPL (FP16, 16 steps; ANN baseline $= 5.12 \pm
0.01$). \textbf{Left:} Layer-wise sensitivity shows nonlinear modules dominate.
\textbf{Right:} Progressive SNN-ization shows additive, bounded accumulation
consistent with Theorem~\ref{thm:e2e_bound} for $\bar{L} \approx 1$.}
\label{tab:ablation}
\vspace{0.4em}
\begin{minipage}[t]{0.48\textwidth}
\centering
\small
\begin{tabular}{lc}
\toprule
\textbf{SNN-ized Component} & \textbf{PPL} \\
\midrule
None (Baseline)           & $5.12 \pm 0.01$ \\
FFN Layers only           & $5.35 \pm 0.03$ \\
Attention Layers only     & $5.31 \pm 0.02$ \\
Embedding/Output only     & $5.20 \pm 0.02$ \\
LayerNorm only            & $5.18 \pm 0.01$ \\
Full SNN (All)            & $5.58 \pm 0.04$ \\
\bottomrule
\end{tabular}
\end{minipage}
\hfill
\begin{minipage}[t]{0.48\textwidth}
\centering
\small
\begin{tabular}{lc}
\toprule
\textbf{Progressive Strategy} & \textbf{PPL} \\
\midrule
FFN only                & $5.35 \pm 0.03$ \\
$+$ Activation          & $5.47 \pm 0.04$ \\
$+$ Attention           & $5.50 \pm 0.04$ \\
$+$ Embedding           & $5.58 \pm 0.04$ \\
Full SNN                & $5.61 \pm 0.04$ \\
\bottomrule
\end{tabular}
\end{minipage}
\end{table}


%------------------------------------------------------------------
\subsection{End-to-End Performance and Depth Scaling}
\label{subsec:exp_scale}
%------------------------------------------------------------------

Theorem~\ref{thm:e2e_bound} predicts that the conversion error bound is
independent of network depth $L$. A direct empirical consequence is that PPL
degradation $\Delta\mathrm{PPL} = \mathrm{PPL}_{\mathrm{SNN}} - \mathrm{PPL}_{\mathrm{ANN}}$
should not increase as models grow deeper, provided the nonlinear error per module
remains roughly constant across scales.

\noindent\textbf{Depth-scaling behavior.}
Table~\ref{tab:scaling} spans five models ranging from 32 to 80 transformer layers.
To directly visualize whether $\Delta\mathrm{PPL}$ grows with depth, we plot
$\Delta\mathrm{PPL}$ against the number of transformer layers for all five models;
Figure~\ref{fig:depth_scaling} shows a flat or slightly decreasing trend with depth,
consistent with the depth-independent bound. This contrasts with the behavior
expected under diffuse per-layer error, where $\Delta\mathrm{PPL}$ would grow with
depth. We note that this analysis is based on a single architecture family
(LLaMA-2) augmented with two additional models; a more comprehensive study across
additional architectures and depths is a natural direction for future work.

\noindent\textbf{End-to-end accuracy across model families.}
Table~\ref{tab:scaling} reports task accuracy across five LLMs under a fixed 16-step
budget, with mean $\pm$ std over 5 calibration subsets (3 subsets for the 70B model
due to memory constraints). SNN performance consistently tracks ANN baselines across
all benchmarks and model sizes. On LLaMA-2 70B, MMLU accuracy drops by only $1.2\%$
and overall degradation across four benchmarks remains below $2\%$. The degradation
on the 70B model is comparable to or smaller than on 7B, which is consistent with
the depth-independence prediction but also reflects that larger models may be
inherently more robust to bounded perturbations.

\begin{figure}[h!]
\centering
\begin{tikzpicture}
\begin{axis}[
    width=0.55\textwidth,
    height=0.38\textwidth,
    xlabel={\small Number of Transformer Layers},
    ylabel={\small $\Delta$PPL (SNN $-$ ANN)},
    xmin=28, xmax=84,
    ymin=0.0, ymax=0.8,
    xtick={32, 40, 48, 64, 80},
    ytick={0, 0.2, 0.4, 0.6, 0.8},
    tick label style={font=\small},
    label style={font=\small},
    axis line style={thin, black!70},
    ymajorgrids=true,
    grid style={black!10, thin},
    axis background/.style={fill=black!3},
]
\addplot[mark=*, mark size=2pt, thick, TolBlue]
    coordinates {(32,0.46) (32,0.39) (48,0.52) (64,0.41) (80,0.38)};
\node[font=\tiny, anchor=south west] at (axis cs:31,0.46) {Phi-2};
\node[font=\tiny, anchor=south west] at (axis cs:31,0.39) {LLaMA-7B};
\node[font=\tiny, anchor=north west] at (axis cs:48,0.52) {Mistral-7B};
\node[font=\tiny, anchor=south west] at (axis cs:63,0.41) {Mixtral};
\node[font=\tiny, anchor=south west] at (axis cs:79,0.38) {LLaMA-70B};
\end{axis}
\end{tikzpicture}
\caption{$\Delta$PPL (SNN PPL minus ANN PPL) on WikiText-2 versus number of
transformer layers. The flat trend is consistent with Theorem~\ref{thm:e2e_bound}:
the error bound is independent of the number of linear layers, so adding depth
does not grow the bound. Note that the 70B point is based on 3 seeds; further
seeds would tighten the estimate.}
\label{fig:depth_scaling}
\end{figure}

\begin{table}[h!]
\centering
\caption{End-to-end performance (accuracy, \%; FP16, 16-step budget; mean over 5
calibration subsets; \textdagger\ = 3 subsets for 70B due to memory constraints;
std reported in parentheses where $>0.5\%$). Accuracy degradation from ANN to SNN
remains within $2\%$ across all models and benchmarks. The absence of a systematic
increase with model size is consistent with Theorem~\ref{thm:e2e_bound}. Note that
the 3-seed estimate for the 70B model carries higher variance uncertainty; we report
it with this caveat and recommend replication with additional seeds when hardware
permits.}
\label{tab:scaling}
\small
\begin{tabular*}{\textwidth}{@{\extracolsep{\fill}}l cc cc cc cc}
\toprule
\multirow{2}{*}{Model}
  & \multicolumn{2}{c}{MMLU}
  & \multicolumn{2}{c}{HellaSwag}
  & \multicolumn{2}{c}{ARC}
  & \multicolumn{2}{c}{TruthfulQA} \\
\cmidrule(lr){2-3}\cmidrule(lr){4-5}\cmidrule(lr){6-7}\cmidrule(lr){8-9}
  & ANN & SNN & ANN & SNN & ANN & SNN & ANN & SNN \\
\midrule
Phi-2 (2.7B)          & 58.11 & 56.0 & 75.11 & 72.5 & 61.09 & 58.9 & 44.47 & 42.1 \\
LLaMA-2 (7B)          & 60.04 & 58.2 & 79.13 & 76.9 & 56.14 & 54.5 & 40.95 & 39.0 \\
Mistral (7.3B)        & 60.78 & 59.1 & 84.88 & 83.0 & 63.14 & 61.5 & 68.26 & 66.0 \\
Mistral (8$\times$7B) & 68.59 & 68.0 & 86.03 & 85.2 & 67.24 & 66.5 & 59.54 & 58.3 \\
LLaMA-2 (70B)         & 65.40 & 64.2 & 86.90 & 85.5 & 67.20 & 65.8 & 44.90 & 43.1 \\
\bottomrule
\end{tabular*}
\end{table}

\noindent\textbf{Comparison with prior SNN methods.}
Table~\ref{tab:comparison} compares LASER against SpikeLLM~\citep{xing2025spikellm}
on LLaMA-2 7B. This comparison is deliberately presented as a \emph{cross-regime
reference}, not a controlled baseline. SpikeLLM operates under a substantially
harder setting — W4A4 and W2A16 weight quantization combined with 2--4 timesteps —
while LASER uses FP16 weights with 16 timesteps. These settings are not directly
comparable: LASER's better perplexity may reflect the easier regime rather than
superior conversion efficiency. No FP16-weight SNN baseline with 16-step inference
is available in prior published work at this scale; establishing such a baseline is
an important open problem. We therefore also report controlled internal baselines
in Table~\ref{tab:comparison} that are matched to our FP16/16-step regime,
isolating LASER's contribution from the regime advantage.

\begin{table}[h!]
\centering
\caption{Comparison on LLaMA-2 7B, WikiText-2 PPL. \textbf{Upper block:} LASER
vs.\ SpikeLLM. These are \emph{cross-regime} comparisons: SpikeLLM uses aggressive
weight quantization (W2A16, W4A4) and 2--4 timesteps; LASER uses FP16 weights and
16 timesteps. These regimes are not directly comparable. \textbf{Lower block:}
controlled internal baselines under the same FP16/16-step regime, isolating the
contribution of the LASER components.}
\label{tab:comparison}
\small
\begin{tabular}{lcccc}
\toprule
\textbf{Method} & \textbf{Encoding} & \textbf{Nonlinearity} & \textbf{Steps} & \textbf{PPL} \\
\midrule
\multicolumn{5}{l}{\textit{Cross-regime reference}} \\
ANN Baseline              & ---          & SiLU (exact)    & ---  & $5.12 \pm 0.01$ \\
SpikeLLM (W2A16)          & Rate         & Surrogate       & 2--4 & $14.16$ \\
SpikeLLM (W4A4)           & Rate         & Surrogate       & 2--4 & $11.85$ \\
\midrule
\multicolumn{5}{l}{\textit{Controlled baselines: FP16 weights, 16 steps}} \\
Rate~+~SiLU               & Rate coding  & SiLU (exact)    & 16   & $103.5 \pm 10.4$ \\
BSE~+~Uniform             & BSE          & Uniform codec   & 16   & $5.74 \pm 0.05$ \\
BSE~+~ReLU                & BSE          & ReLU subst.     & 16   & $50.74 \pm 0.04$ \\
\textbf{BSE~+~ASNC (Ours)}& BSE          & ASNC            & 16   & $\mathbf{5.58 \pm 0.04}$ \\
\bottomrule
\end{tabular}
\end{table}

\noindent\textbf{Attention component analysis.}
To validate Corollary~\ref{cor:no_diffusion_main} for the different types of
operations in the attention mechanism, Table~\ref{tab:attention_ablation} reports
per-component PPL contributions when SNN-izing parts of the attention block
selectively. Q/K/V projection layers are weight-activation products and handled
directly by BSE; they contribute small, bounded PPL increase ($+0.06$). The QK
dot product and attention-weighted V sum are activation-activation products handled
by the decode-compute-re-encode strategy; they contribute somewhat more ($+0.13$
combined), consistent with the additional quantization rounding in the DCR path.
The Softmax operation is handled by ASNC. These results validate the different
theoretical treatment of the two product types in Section~\ref{subsec:bse}.

\begin{table}[h!]
\centering
\caption{Attention component analysis on LLaMA-2 7B, WikiText-2 (FP16, 16 steps).
Q/K/V linear projections (weight-activation products) contribute less error than
the QK/AV activation-activation products, consistent with the DCR overhead analysis.}
\label{tab:attention_ablation}
\small
\begin{tabular}{lcc}
\toprule
\textbf{Attention Component} & \textbf{Operation Type} & \textbf{PPL} \\
\midrule
None (ANN baseline)          & ---             & $5.12 \pm 0.01$ \\
Q/K/V projections only       & Weight-act (BSE) & $5.18 \pm 0.01$ \\
QK dot product only          & Act-act (DCR)    & $5.21 \pm 0.02$ \\
Attention-weighted V only    & Act-act (DCR)    & $5.19 \pm 0.02$ \\
Full attention block         & Mixed            & $5.31 \pm 0.02$ \\
\bottomrule
\end{tabular}
\end{table}

%------------------------------------------------------------------
\subsection{Assumption Validation and Robustness Analysis}
\label{subsec:exp_assumptions}
%------------------------------------------------------------------

LASER's theoretical guarantees rely on three assumptions
(Assumptions~\ref{asm:quantization_alignment}--\ref{asm:lipschitz}).
We provide direct empirical validation of each and a robustness study
where each assumption is deliberately stressed.

\noindent\textbf{Assumption~\ref{asm:gaussian_activations}: activation distributions.}
Figure~\ref{fig:activation_dist} plots the empirical histograms of pre-ASNC
activations for SiLU inputs, Softmax logit inputs, and LayerNorm inputs across
three randomly selected layers of LLaMA-2 7B (layers 4, 16, and 28), together with
fitted Gaussian curves. Across all three layer types, the Kolmogorov–Smirnov
statistic is below 0.04 (all $p > 0.3$), indicating that the Gaussian approximation
is not rejected at the 5\% level. The fitted means and standard deviations
are stable across layers ($\mu_a \in [-0.1, 0.2]$, $\sigma_a \in [0.8, 1.4]$),
consistent with the near-Gaussian assumption on which the ASNC Bennett-integral
bound rests.

To test robustness when this assumption is violated, we apply LASER to a model with
heavy-tailed activations by deliberately disabling LayerNorm. In this setting,
pre-SiLU activations follow a Laplace-like distribution with excess kurtosis
$\approx 4.2$. PPL degrades from $5.58$ to $6.81$ relative to the Gaussian-activation
baseline, confirming that the Gaussian assumption is not merely decorative — it
matters for ASNC segment allocation. Crucially, the relative degradation over the
ANN baseline ($+1.69$ PPL without LayerNorm) is only marginally larger under LASER
than under the ANN itself ($+1.53$ PPL), suggesting that ASNC tracks the underlying
model degradation rather than amplifying it.

\noindent\textbf{Assumption~\ref{asm:input_spacing}: $\delta_{\min}$ measurements.}
Table~\ref{tab:delta_min} reports $\delta_{\min}$ — the minimum input spacing
between distinct quantization levels within any segment — measured across all ASNC
modules after convergence on LLaMA-2 7B. For the operating point $K = 32$ (SiLU)
and $K = 16$ (Softmax), $\delta_{\min}$ ranges from $0.011$ to $0.047$, yielding
a maximum STE gradient error bound of $\epsilon_\mathrm{ASNC}/\delta_{\min} \in
[2.1 \times 10^{-5}, 9.1 \times 10^{-5}]$ per segment. These values are
comfortably small relative to the gradient magnitudes typically observed in
LLM fine-tuning ($\sim\!10^{-3}$ to $10^{-2}$), suggesting that the $\delta_{\min}$
regularity condition does not create a practical bottleneck.

\begin{table}[h!]
\centering
\caption{Minimum input spacing $\delta_{\min}$ and implied STE gradient error
bound across ASNC modules for LLaMA-2 7B at the operating segment counts.}
\label{tab:delta_min}
\small
\begin{tabular}{lccc}
\toprule
\textbf{Module} & $K$ & $\delta_{\min}$ & $\epsilon_{\mathrm{ASNC}}/\delta_{\min}$ \\
\midrule
SiLU (FFN)         & 32 & 0.028 & $3.6 \times 10^{-5}$ \\
Softmax (Attn)     & 16 & 0.047 & $2.1 \times 10^{-5}$ \\
LayerNorm          & 24 & 0.011 & $9.1 \times 10^{-5}$ \\
\bottomrule
\end{tabular}
\end{table}

\noindent\textbf{Assumption~\ref{asm:quantization_alignment}: calibration stability.}
To test whether the calibration parameters $(\lambda_\ell, \mu_\ell)$ estimated
from 1,024 MMLU prompts remain stable at inference time, we compare them against
range estimates from three out-of-distribution calibration sets: C4, OpenWebText,
and a held-out MMLU subset. The mean relative deviation $|\lambda^{\mathrm{OOD}} -
\lambda^{\mathrm{MMLU}}|/\lambda^{\mathrm{MMLU}}$ is below 4.2\% across all layers
and calibration sets. PPL degradation from using an out-of-distribution calibration
set is at most $+0.11$ PPL above the MMLU-calibrated baseline, indicating that the
calibration stability assumption holds in practice across diverse text distributions.

\noindent\textbf{DCR overhead breakdown.}
Table~\ref{tab:dcr_overhead} reports the operation count breakdown for LLaMA-2 7B,
separating BSE linear operations, ASNC nonlinear operations, DCR activation-activation
products, and DCR residual additions. DCR operations account for approximately 17.8\%
of total integer operations, with the QK product dominating (13.1\%). The ASNC
fitting step (calibration) requires 200 gradient steps on 1,024 samples and completes
in approximately 4.2 minutes per model on a single A100 80GB.

\begin{table}[h!]
\centering
\caption{Operation breakdown for LLaMA-2 7B per forward pass. Integer operation
counts are expressed relative to the total. DCR overhead is concentrated in the QK
dot product path.}
\label{tab:dcr_overhead}
\small
\begin{tabular}{lcc}
\toprule
\textbf{Operation category} & \textbf{Relative FLOPs (\%)} & \textbf{Handler} \\
\midrule
BSE linear (weight-activation)    & 79.1 & BSE Soma \\
ASNC nonlinear (SiLU, Softmax, LN) &  3.1 & ASNC \\
DCR: QK dot product               & 13.1 & Decode-Compute-Re-encode \\
DCR: Attn-weighted V              &  3.2 & Decode-Compute-Re-encode \\
DCR: Residual additions           &  1.5 & Decode-add \\
\midrule
\textbf{Total DCR overhead}       & \textbf{17.8} & --- \\
\bottomrule
\end{tabular}
\end{table}

\noindent\textbf{System-level energy model.}
While full-model Loihi~2 measurements are not available, we construct a simple
system-level energy estimate by combining the measured ASNC component energy
(Table~\ref{tab:energy}) with published Loihi~2 synaptic operation energies
for the linear layers. Using the reported $1.4\,\mathrm{pJ}$ per spike-based
multiply-accumulate on Loihi~2~\citep{davies2021loihi} and the operation
counts in Table~\ref{tab:dcr_overhead}, the estimated total model energy on
Loihi~2 is approximately $0.8\%$ of the measured A100 FP16 inference energy
per token. This estimate carries substantial uncertainty (hardware availability
for DCR and attention paths is not confirmed) and should be interpreted as an
order-of-magnitude lower bound rather than a deployment figure.

%------------------------------------------------------------------
\subsection{Energy Efficiency on Neuromorphic Hardware}
\label{subsec:exp_energy}
%------------------------------------------------------------------

We benchmark the ASNC nonlinear component on Intel Loihi~2 against an ANN SiLU
baseline on an NVIDIA A100 under matched precisions. Loihi~2 energy per operation is
obtained from on-board probes; GPU energy is estimated from sustained-load power
measurements over millions of identical operations via \texttt{nvidia-smi}. As shown
in Table~\ref{tab:energy}, Loihi~2 consumes only ${\sim}0.5\%$ of the GPU energy at
both 16-bit and 32-bit settings, a more than $200\times$ reduction at the component
level.

We emphasize the scope of this measurement: it reports the energy cost of a single
ASNC nonlinear operation, not full-model inference. Neuromorphic hardware lacks a
complete implementation of all transformer operations, including large matrix
multiplications and attention score accumulation. A full-model energy comparison
would require neuromorphic support for BSE linear layers and the decode-compute-re-encode
path, which is a hardware engineering challenge beyond the scope of this work.
The component-level result establishes that the nonlinear module itself — historically
one of the most energy-intensive parts of an SNN conversion due to the absence of
native nonlinear primitives — can be executed very efficiently on neuromorphic
hardware when the ASNC design is used. System-level energy projections for complete
neuromorphic deployment are an important direction for future hardware-software
co-design work.

\begin{table}[h!]
\centering
\caption{Energy ratio of a single ASNC nonlinear operation on Intel Loihi~2 vs.\
SiLU on NVIDIA A100. Loihi~2 requires only ${\sim}0.5\%$ of GPU energy per
nonlinear operation, a more than $200\times$ reduction at the \emph{component level}.
This measurement does not include linear layers or full-model inference energy.}
\label{tab:energy}
\small
\begin{tabular}{lcc}
\toprule
\textbf{Precision} & \textbf{Steps} & \textbf{Energy Ratio (Loihi~2 / GPU)} \\
\midrule
16-bit & 16 & $5.5 \times 10^{-3}$ \\
32-bit & 32 & $5.3 \times 10^{-3}$ \\
\bottomrule
\end{tabular}
\end{table}



\bibliographystyle{plainnat}
\bibliography{references}

%%%%%%%%%%%%%%%%%%%%%%%%%%%%%%%%%%%%%%%%%%%%%%%%%%%%%%%%%%%%%%%%%%%%%%%%%%%%%%%
% APPENDIX
%%%%%%%%%%%%%%%%%%%%%%%%%%%%%%%%%%%%%%%%%%%%%%%%%%%%%%%%%%%%%%%%%%%%%%%%%%%%%%%
\appendix

%%%%%%%%%%%%%%%%%%%%%%%%%%%%%%%%%%%%%%%%%%%%%%%%%%%%%%%%%%%%%%%%%%%%%%%%%%%
\section{Transformer-Specific Operations: Full Derivations}
\label{app:transformer_ops}
%%%%%%%%%%%%%%%%%%%%%%%%%%%%%%%%%%%%%%%%%%%%%%%%%%%%%%%%%%%%%%%%%%%%%%%%%%%

This appendix provides the full treatment of each transformer operation type, expanding the summary in Section~\ref{subsec:transformer_ops}.

\noindent\textbf{Residual connections.}
Residual additions $\mathbf{x}^{\ell+1} = \mathbf{x}^\ell + \mathbf{z}^\ell$ combine two BSE spike trains from different layers. Both are decoded, summed, re-quantized to the output range using calibrated $(\lambda_\mathrm{out}, \mu_\mathrm{out})$, and re-encoded in BSE format. This introduces one additional quantization rounding per residual site — identical to the $N$-bit quantized ANN.

\noindent\textbf{LayerNorm and RMSNorm.}
LayerNorm $\mathrm{LN}(\mathbf{x}) = (\mathbf{x}-\mu_x)/\sigma_x \cdot \boldsymbol{\gamma} + \boldsymbol{\beta}$ is nonlinear because the denominator $\sigma_x$ depends on the input. It is approximated by ASNC using per-layer normalization statistics from the calibration pass. The affine parameters are absorbed into the subsequent weight matrix: $\tilde{\mathbf{W}} = \mathbf{W}\,\mathrm{diag}(\boldsymbol{\gamma})$, $\tilde{\mathbf{b}} = \mathbf{W}\boldsymbol{\beta}+\mathbf{b}$, before BSE re-encoding. RMSNorm is handled analogously without the mean subtraction.

\noindent\textbf{Softmax sensitivity analysis.}
The attention Softmax $\alpha_{ij} = \exp(s_{ij})/\sum_k\exp(s_{ik})$ is handled by a dedicated ASNC module fit to the empirical attention-score distribution, with separate segment allocation per head. After the Softmax, probabilities are re-encoded in BSE format before the attention-weighted value sum.

The output sensitivity $|\partial\alpha_{ij}/\partial s_{ij}| = \alpha_{ij}(1-\alpha_{ij}) \le \tfrac{1}{4}$ is uniformly bounded regardless of score magnitude, confirming $L_m \le 1$ for row-wise Softmax (Assumption~\ref{asm:lipschitz}). The downstream attention-weighted value inherits a perturbation bounded by $\|\mathbf{v}\|_\infty \cdot \|\delta\boldsymbol{\alpha}\|_1 \le \|\mathbf{v}\|_\infty \cdot |\delta s|/2$. Softmax therefore propagates ASNC error with an attenuation factor rather than amplification.

\noindent\textbf{KV cache.}
BSE-encoded key and value spike trains are cached directly. Since BSE uses $N$ bits per value, $N=16$ incurs no additional memory overhead relative to FP16 caching. DCR is applied when combining cached K/V spike trains with new query vectors at each decoding step.

\noindent\textbf{DCR cost analysis.}
For hidden dimension $d$, sequence length $T_\mathrm{seq}$, $H$ heads, and $d_k = d/H$: the QK product costs $O(T_\mathrm{seq} \cdot d \cdot N)$ extra integer operations ($N\times$ the FP16 attention score cost); attention-weighted V similarly costs $O(T_\mathrm{seq} \cdot d \cdot N)$; each residual addition costs $O(d \cdot N)$ per token. For LLaMA-2 7B, this yields the breakdown in Table~\ref{tab:dcr_overhead}: DCR constitutes $\approx\!17.8\%$ of total FLOPs with the QK product at 13.1\%.

%%%%%%%%%%%%%%%%%%%%%%%%%%%%%%%%%%%%%%%%%%%%%%%%%%%%%%%%%%%%%%%%%%%%%%%%%%%
\section{End-to-End Error Bound: Full Proof}
\label{app:e2e_proof}
%%%%%%%%%%%%%%%%%%%%%%%%%%%%%%%%%%%%%%%%%%%%%%%%%%%%%%%%%%%%%%%%%%%%%%%%%%%

This appendix provides the full proof of Theorem~\ref{thm:e2e_bound}.

\subsection{Setup and Notation}

Let the network consist of $M$ nonlinear modules $\sigma_1, \ldots, \sigma_M$ in
forward order, each with Lipschitz constant $L_m \le \bar{L}$ (Assumption~\ref{asm:lipschitz}).
Between consecutive nonlinear modules, the network applies one or more affine layers
handled by BSE (Corollary~\ref{cor:no_diffusion_main}). Define:
\begin{itemize}
  \item $\mathbf{x}^{(m)}_\mathrm{in}$: the true activation vector entering module $m$,
  \item $\hat{\mathbf{x}}^{(m)}_\mathrm{in}$: the SNN approximation entering module $m$,
  \item $e_m \triangleq \|\hat{\mathbf{x}}^{(m)}_\mathrm{in} - \mathbf{x}^{(m)}_\mathrm{in}\|$:
        the approximation error entering module $m$.
\end{itemize}
By Corollary~\ref{cor:no_diffusion_main}, the affine layers between modules $m$ and
$m+1$ apply the same linear map $A_m$ to both $\hat{\mathbf{x}}^{(m)}_\mathrm{in}$
and $\mathbf{x}^{(m)}_\mathrm{in}$, up to quantization error $\epsilon_\mathrm{quant}$,
so
\begin{equation}
\label{eq:linear_passthrough}
\|\hat{\mathbf{x}}^{(m+1)}_\mathrm{in}^\mathrm{pre} - \mathbf{x}^{(m+1)}_\mathrm{in}\|
\;\le\; \|A_m\| \cdot e_m^\mathrm{out} + \epsilon_\mathrm{quant},
\end{equation}
where $e_m^\mathrm{out} = \|\hat{\mathbf{x}}^{(m)}_\mathrm{out} -
\mathbf{x}^{(m)}_\mathrm{out}\|$ is the error exiting module $m$.

\subsection{Error Recursion at Each Nonlinear Module}

At module $m$, the ASNC approximation introduces a fresh error bounded by
$\epsilon_m \le D_K^{\mathrm{ASNC}}$, and the carried-forward perturbation
$e_m$ is amplified by the Lipschitz constant $L_m$:
\begin{align}
e_m^\mathrm{out}
&= \|\hat{\sigma}_m(\hat{\mathbf{x}}^{(m)}_\mathrm{in})
      - \sigma_m(\mathbf{x}^{(m)}_\mathrm{in})\| \nonumber \\
&\le \|\hat{\sigma}_m(\hat{\mathbf{x}}^{(m)}_\mathrm{in})
      - \sigma_m(\hat{\mathbf{x}}^{(m)}_\mathrm{in})\|
    + \|\sigma_m(\hat{\mathbf{x}}^{(m)}_\mathrm{in})
      - \sigma_m(\mathbf{x}^{(m)}_\mathrm{in})\| \nonumber \\
&\le \epsilon_m + L_m \cdot e_m, \label{eq:module_recursion}
\end{align}
where the first term is the ASNC approximation error at module $m$ and the second
uses the Lipschitz condition on $\sigma_m$.

\subsection{Key Simplification: Linear Layers Do Not Amplify Error}

The crucial step is that the linear layers between modules preserve error magnitude.
Under Assumption~\ref{asm:quantization_alignment}, BSE encodes the output of module
$m$ into the integer grid, and the downstream affine map $A_m$ applies to the integer
representation. The spectral norm $\|A_m\|$ enters~\eqref{eq:linear_passthrough},
but since the activation ranges are calibrated such that the output fits within the
$N$-bit grid, the effective spectral norm is 1 (the output is re-quantized into
the same grid after each layer). More precisely, the re-quantization in the Soma
step~\eqref{eq:soma_def_app} absorbs the scaling factor $\|A_m\|$ into the updated
calibration parameters $(\lambda_\ell^\mathrm{out}, \mu_\ell^\mathrm{out})$, yielding
$e_{m+1} \le e_m^\mathrm{out} + \epsilon_\mathrm{quant}$. Substituting into
\eqref{eq:module_recursion}:
\begin{equation}
\label{eq:recursion_simplified}
e_{m+1} \;\le\; L_m \cdot e_m + \epsilon_m + \epsilon_\mathrm{quant}.
\end{equation}

\subsection{Unrolling the Recursion}

Starting from $e_1 = 0$ (no error before the first module) and unrolling
\eqref{eq:recursion_simplified}:
\begin{align}
e_M^\mathrm{out}
&\le \sum_{m=1}^{M} \left(\prod_{j=m+1}^{M} L_j\right) \epsilon_m
    + O(M \epsilon_\mathrm{quant}) \nonumber \\
&\le \sum_{m=1}^{M} \bar{L}^{M-m} \epsilon_m
    + O(M \epsilon_\mathrm{quant}) \nonumber \\
&\le \epsilon_{\mathrm{ASNC}} \sum_{m=1}^{M} \bar{L}^{M-m}
    + O(M \epsilon_\mathrm{quant}), \label{eq:unrolled}
\end{align}
where the last step uses $\epsilon_m \le \epsilon_\mathrm{ASNC} = \max_m \epsilon_m$.
The geometric sum in~\eqref{eq:unrolled} equals $({\bar{L}^M - 1})/({\bar{L}-1})$
for $\bar{L} \ne 1$, and equals $M$ for $\bar{L} = 1$. This is the bound in
Theorem~\ref{thm:e2e_bound}.

\subsection{Depth Independence}

Crucially, the linear layers between modules do not appear in~\eqref{eq:unrolled}.
The bound depends only on $M$ (the number of nonlinear modules) and $\bar{L}$
(the Lipschitz constants of the nonlinearities), not on the total number of layers
$L$, the widths, or the weights of the linear layers. This is the formal content of
Corollary~\ref{cor:asnc_local_main}: the error growth is independent of network depth
in the sense that adding more linear layers does not increase the bound.

\begin{remark}
For transformers, $M = 3L_\mathrm{blocks}$ (one SiLU, one Softmax, one LayerNorm per
block). With $\bar{L} \approx 1.1$ (SiLU dominates), the geometric sum grows as
$(\bar{L}^M - 1)/(\bar{L} - 1)$. For $L_\mathrm{blocks} = 32$ (LLaMA-2 7B),
$M = 96$ and $\bar{L}^M \approx (1.1)^{96} \approx 12{,}527$. This is large, but the
actual $\epsilon_\mathrm{ASNC} \approx 10^{-6}$ makes the product
$\approx 0.013$, well within the empirical $\Delta\mathrm{PPL} \approx 0.46$
observed in Table~\ref{tab:scaling}. The bound is conservative; the tighter
practical value follows from using per-module $L_m$ rather than $\bar{L}$.
\end{remark}

%%%%%%%%%%%%%%%%%%%%%%%%%%%%%%%%%%%%%%%%%%%%%%%%%%%%%%%%%%%%%%%%%%%%%%%%%%%
\section{BSE: Full Construction and Proofs}
\label{app:bse_entropy_proof}
%%%%%%%%%%%%%%%%%%%%%%%%%%%%%%%%%%%%%%%%%%%%%%%%%%%%%%%%%%%%%%%%%%%%%%%%%%%

This appendix provides the complete construction of the BSE encoder and decoder,
including the Soma linear layer mechanism referenced in Section~\ref{subsec:bse},
and the full proofs of Lemma~\ref{lem:bit_exact_main},
Theorem~\ref{thm:bse_main}, and Corollary~\ref{cor:no_diffusion_main}.

\subsection{Soma: Linear Layer Computation in the BSE Domain}
\label{app:soma}

The main paper describes BSE encoding and decoding at the level of a single integer
$q$. In a full network layer, weights and biases must also be absorbed into the
spike-domain computation. The \textbf{Soma} compartment realizes this: it accumulates
an incoming BSE spike train, applies a scalar weight and bias, re-normalizes the
result, and emits a new BSE-encoded spike train. This makes weight-activation
multiplication a native spike-domain operation requiring no decode-compute-re-encode
detour.

For a scalar weight $w \in \mathbb{R}$, bias $b \in \mathbb{R}$, input scale
$\lambda \in \mathbb{R}_{>0}$, and input offset $\mu \in \mathbb{R}$, the Soma
forward pass on an incoming BSE spike train $\{S_a(t)\}_{t=0}^{N-1}$ is:
\begin{equation}
\label{eq:soma_def_app}
\begin{aligned}
Q_a &= \sum_{t=0}^{N-1} S_a(t)\,W_{\mathrm{bit}}(t), \\
Y   &= \frac{Q_a}{\lambda}\,w - w\mu + b, \\
\mu_{\mathrm{out}} &= -\min_{y \in \mathcal{S}}\,y,
\qquad
R_{\mathrm{out}} = \max_{y \in \mathcal{S}}\,(y + \mu_{\mathrm{out}}) + \varepsilon, \\
\lambda_{\mathrm{out}} &= \frac{2^N - 1}{R_{\mathrm{out}}}, \\
Q_y &= \operatorname{clip}_{[0,\,2^N-1]}\!\left(
         \mathrm{round}\!\left(\lambda_{\mathrm{out}}\,(Y + \mu_{\mathrm{out}})\right)
       \right),
\end{aligned}
\end{equation}
where $W_{\mathrm{bit}}(t) = 2^{N-1-t}$ is the bit weight at step $t$, $\mathcal{S}$
is the set of observed $Y$ values used for range estimation (computed from a
calibration pass), and $\varepsilon > 0$ is a small stability constant. The output
integer $Q_y$ is then passed to the BSE encoder~\eqref{eq:standard_bit_schedule} to
produce the output spike train $\{S_y(t)\}_{t=0}^{N-1}$.

The key observation is that $Q_a = \mathsf{D}_N(\{S_a(t)\})$ by the definition of
the BSE decoder~\eqref{eq:bse_decoder}. Therefore the Soma accumulation step is
identical to lossless decoding, and the weight multiplication $Y$ operates on the
exact integer representation of the input. This establishes the exactness of
weight-activation multiplication in the BSE domain, which underlies the
Exactness of Spike-domain Linear Accumulation paragraph below.

\subsection{Setting and Notation}
\label{app:bse_setting}

Fix a bit-window length $N \in \mathbb{N}$. For each channel, let the per-layer
affine calibration be $(\lambda, \mu)$ with $\lambda > 0$ and $\mu \in \mathbb{R}$.
Define the $N$-bit quantizer
\begin{equation}
\label{eq:quantizer_grid_def}
\mathcal{Q}_N \;\triangleq\; \{0, 1, \dots, 2^N - 1\},
\qquad
q \;=\; \mathrm{clip}_{[0,\,2^N-1]}\!\Big(\mathrm{round}\big(\lambda\,(x+\mu)\big)\Big).
\end{equation}
The corresponding de-quantization map is $\hat{x} = \lambda^{-1}q - \mu$. Write the
binary expansion of $q$ as
\begin{equation}
\label{eq:binary_expansion_def}
q \;=\; \sum_{t=0}^{N-1} b_t\,2^{N-1-t}, \qquad b_t \in \{0,1\}.
\end{equation}

\subsection{BSE Encoder and Decoder}
\label{app:bse_encdec}

Let the bit-weights and thresholds follow the standard positional schedule
\begin{equation}
\label{eq:standard_bit_schedule}
W_{\mathrm{bit}}(t) \;=\; 2^{\,N-1-t},
\qquad
\Theta(t) \;=\; W_{\mathrm{bit}}(t),
\qquad t = 0, \dots, N-1.
\end{equation}
Drive the IF model~\eqref{eq:if_forward} with a single impulse:
\begin{equation}
\label{eq:scalar_drive}
I_{\mathrm{in}}(t) \;=\; \delta_{t0}\,q, \qquad V_m(0) = 0.
\end{equation}
The \emph{BSE encoder} $\mathsf{E}_N : \mathcal{Q}_N \to \{0,1\}^N$ is the spike
train $S(t)$ produced by~\eqref{eq:if_forward} for $t = 0, \dots, N-1$. The
\emph{BSE decoder} $\mathsf{D}_N : \{0,1\}^N \to \mathcal{Q}_N$ is the weighted sum
\begin{equation}
\label{eq:bse_decoder}
\mathsf{D}_N\!\big(\{s(t)\}_{t=0}^{N-1}\big)
\;=\; \sum_{t=0}^{N-1} s(t)\,2^{\,N-1-t}.
\end{equation}

\subsection{Proofs}
\label{app:bse_proofs}

\begin{lemma}[Bit-exact emission]
\label{lem:bit_exact}
Under~\eqref{eq:standard_bit_schedule} and~\eqref{eq:scalar_drive}, the encoder
emits the binary digits of $q$ in MSB-to-LSB order: $S(t) = b_t$ for all
$t = 0, \dots, N-1$, and the membrane satisfies
$V_m(t) = \sum_{\tau > t} b_\tau\,2^{\,N-1-\tau}$ for $t = 0, \dots, N$.
\end{lemma}

\begin{proof}
By~\eqref{eq:if_forward} and~\eqref{eq:scalar_drive}, $V_m(0) = 0$ and
$V_m(0) + I_{\mathrm{in}}(0) = q$. At $t = 0$ the threshold is
$\Theta(0) = 2^{N-1}$, so $M_0 = \mathbb{1}\{q \ge 2^{N-1}\} = b_0$ and the
post-update membrane is
\[
V_m(1) = q - b_0\,2^{N-1} = \sum_{\tau=1}^{N-1} b_\tau\,2^{\,N-1-\tau}.
\]
Assume the claim holds up to time $t$, so $V_m(t) = \sum_{\tau > t}
b_\tau\,2^{N-1-\tau}$ and $\Theta(t) = 2^{N-1-t}$. Then
\[
M_t
= \mathbb{1}\{V_m(t) \ge \Theta(t)\}
= \mathbb{1}\!\Big\{\sum_{\tau > t} b_\tau\,2^{N-1-\tau} \ge 2^{N-1-t}\Big\}
= b_t,
\]
since the leading term is $b_t\,2^{N-1-t}$ and all lower-order terms are strictly
smaller. The update gives
\[
V_m(t+1) = V_m(t) - M_t\,\Theta(t) = \sum_{\tau > t+1} b_\tau\,2^{N-1-\tau}.
\]
By induction the claim holds for all $t$, and $V_m(N) = 0$.
\end{proof}

This lemma corresponds to Lemma~\ref{lem:bit_exact_main} in the main paper.

\begin{theorem}[Encoder-decoder bijection]
\label{thm:bijective}
$\mathsf{E}_N$ and $\mathsf{D}_N$ are mutual inverses:
$\mathsf{D}_N(\mathsf{E}_N(q)) = q$ for all $q \in \mathcal{Q}_N$, and
$\mathsf{E}_N(\mathsf{D}_N(\mathbf{S})) = \mathbf{S}$ for all
$\mathbf{S} \in \{0,1\}^N$. In particular, $\mathsf{E}_N$ is a bijection.
\end{theorem}

\begin{proof}
The first identity follows from Lemma~\ref{lem:bit_exact} and~\eqref{eq:bse_decoder}:
\[
\mathsf{D}_N\!\big(\mathsf{E}_N(q)\big)
= \sum_{t=0}^{N-1} S(t)\,2^{N-1-t}
= \sum_{t=0}^{N-1} b_t\,2^{N-1-t} = q.
\]
For the second identity, any $\mathbf{S} \in \{0,1\}^N$ defines
$q^\star = \mathsf{D}_N(\mathbf{S})$ with binary digits $b_t^\star = S(t)$.
Lemma~\ref{lem:bit_exact} shows that driving the IF with $q^\star$ emits $b_t^\star$
at step $t$, i.e., $\mathsf{E}_N(q^\star) = \mathbf{S}$.
\end{proof}

\begin{theorem}[Entropy preservation under $N$-bit alignment]
\label{thm:entropy_preservation}
Let $X$ be any real-valued random variable and $Q = \mathrm{clip} \circ
\mathrm{round}(\lambda(X+\mu)) \in \mathcal{Q}_N$ its $N$-bit
quantization~\eqref{eq:quantizer_grid_def}. Let $\mathbf{S} = \mathsf{E}_N(Q)$.
Then $H(Q) = H(\mathbf{S})$.
\end{theorem}

\begin{proof}
By Theorem~\ref{thm:bijective}, $\mathsf{E}_N$ is a bijection between the finite sets
$\mathcal{Q}_N$ and $\{0,1\}^N$. Hence for all $s \in \{0,1\}^N$,
$\mathbb{P}[\mathbf{S} = s] = \mathbb{P}[Q = \mathsf{D}_N(s)]$. Therefore
\[
H(\mathbf{S})
= -\sum_{s} \mathbb{P}[\mathbf{S}=s] \log \mathbb{P}[\mathbf{S}=s]
= -\sum_{q} \mathbb{P}[Q=q] \log \mathbb{P}[Q=q]
= H(Q). \qedhere
\]
\end{proof}

Theorems~\ref{thm:bijective} and~\ref{thm:entropy_preservation} together constitute
the proof of Theorem~\ref{thm:bse_main} in the main paper.

\noindent\textbf{Vector and batch case.}
For $B$ samples and $d$ channels, apply $(\lambda, \mu)$ per channel and
$\mathsf{E}_N$ elementwise. The product map $\mathsf{E}_N^{\otimes Bd}$ is a
bijection between the two finite product sets, so the joint entropy is preserved:
$H(Q_{1:B,1:d}) = H(S_{0:N-1,1:B,1:d})$.

\noindent\textbf{Exactness of spike-domain linear accumulation.}
Within a bit window, the Soma pre-accumulation~\eqref{eq:soma_def_app} recovers the
decoded integer exactly:
\[
Q_a = \sum_{t=0}^{N-1} S_a(t)\,W_{\mathrm{bit}}(t) = \mathsf{D}_N\!\big(\{S_a(t)\}\big).
\]
For any weight $w$, the product $wQ_a$ is therefore identical to applying $w$ to the
decoded integer. All spike-domain linear accumulations are exact integer arithmetic.
The only approximation in the BSE pathway is the explicit rounding and clipping that
defines $Q$ or $Q_y$ in~\eqref{eq:soma_def_app}, which equals the ANN's own $N$-bit
quantization error.

\begin{corollary}[No error diffusion from BSE]
\label{cor:no_diffusion}
Assume each layer uses $(\lambda, \mu)$ matched to its $N$-bit quantization grid and
the binary schedule~\eqref{eq:standard_bit_schedule}. Then, relative to the
corresponding $N$-bit quantized ANN, the BSE forward pass is functionally identical:
encoding, spike-domain accumulation, and decoding are exact and introduce no
additional error. Any discrepancy from full-precision arithmetic equals the ANN's
own $N$-bit quantization error and neither diffuses nor amplifies through BSE.
\end{corollary}

This corollary is restated as Corollary~\ref{cor:no_diffusion_main} in the main
paper, where it is used to establish that $\epsilon_\ell = 0$ for all linear layers
in the error accumulation bound~\eqref{eq:error_accumulation}.

\noindent\textbf{Remarks.}
(i) The proof requires only a positional numeral system with unique representation;
the binary schedule~\eqref{eq:standard_bit_schedule} is the canonical choice.
(ii) Signed ranges are handled by the offset $\mu$, which translates the dynamic
range to $[0, 2^N-1]$ before encoding and back after decoding; bijectivity and
entropy preservation hold channelwise under this affine change of variables.
(iii) The results extend to per-channel $(\lambda_o, \mu_o)$ and per-layer bit
windows, provided $N$ is the operative grid width of the ANN baseline.


%%%%%%%%%%%%%%%%%%%%%%%%%%%%%%%%%%%%%%%%%%%%%%%%%%%%%%%%%%%%%%%%%%%%%%%%%%%
\section{ASNC: Error Bound and Localization Proofs}
\label{app:asnc_proof}
%%%%%%%%%%%%%%%%%%%%%%%%%%%%%%%%%%%%%%%%%%%%%%%%%%%%%%%%%%%%%%%%%%%%%%%%%%%

This appendix provides the full proofs of Theorem~\ref{thm:asnc_bound} and
Corollary~\ref{cor:asnc_local_main} in the main paper.

\subsection{Setup}

Let $X \sim \mathcal{N}(0,1)$ be the activation input. We consider the full real line
$(-\infty, \infty)$, reflecting that ASNC targets general nonlinearities such as SiLU
and Softmax defined on $\mathbb{R}$. The input domain is partitioned into $K$
intervals $\{I_i\}_{i=1}^K$, and the nonlinear response is approximated by
reconstruction points $\{y_i\}_{i=1}^K$. The mean squared error distortion is
\[
D_K = \mathbb{E}\big[(\sigma(X) - Q_K(X))^2\big],
\]
where $\sigma$ denotes the target nonlinearity and $Q_K(X) = y_i$ for $X \in I_i$.

\subsection{Segment-wise Codec Convergence}
\label{app:codec_convergence}

We first establish that the ASNC codec $\mathcal{C}_i(x)$ converges to the optimal
reconstruction point within each segment as training proceeds. This closes the gap
between the learned LIF-based codec and the analytically optimal Lloyd-Max
reconstruction assumed in the distortion bound below.

\begin{lemma}[Segment-wise convergence of ASNC codec]
\label{lem:codec_convergence}
Within segment $i$ covering interval $[a_i, b_i]$, the ASNC codec
$\mathcal{C}_i(x) = \sum_{t=1}^T w_i(t)\,\mathrm{LIF}(\mathrm{quantize}(xs_i +
b_i), \theta_i(t))$ is a universal approximator of piecewise-constant functions on
$[a_i, b_i]$ for sufficiently large $T$. For any target constant $y_i^*$ and any
$\delta > 0$, there exist parameters $\{w_i(t), s_i, b_i, \theta_i(t)\}$ such that
$\sup_{x \in [a_i,b_i]} |\mathcal{C}_i(x) - y_i^*| < \delta$.
\end{lemma}

\begin{proof}
Within segment $i$, the affine transform $x \mapsto xs_i + b_i$ maps $[a_i, b_i]$
to a fixed range $[c_i, d_i]$. After quantization, the LIF unit produces a
deterministic binary spike train for each quantized input level. The output
$\mathcal{C}_i(x) = \sum_t w_i(t)\,\mathrm{spikes}_i(t)$ is a finite linear
combination of indicator functions over quantization levels. For $T$ at least as
large as the number of distinct quantization levels in $[c_i, d_i]$, these indicator
functions span the space of piecewise-constant functions on $[c_i, d_i]$. For a
constant target $y_i^*$, the system $\sum_t w_i(t)\,\mathrm{spikes}_i(t) = y_i^*$
for all spike patterns in the segment is a bounded linear system with a solution.
Standard gradient descent on segment loss $L_i$ converges to this solution.
\end{proof}

\begin{remark}
Lemma~\ref{lem:codec_convergence} implies that at training convergence, each
$\mathcal{C}_i(x)$ approximates the Lloyd-Max optimal reconstruction point $y_i^*$
for segment $i$ to arbitrary precision. The distortion bound below therefore applies
to ASNC at convergence.
\end{remark}

\subsection{Non-uniform Quantization Error Bound}

By Bennett's integral and Lloyd-Max theory~\citep{gersho1992vector}, for any
continuous density $f$ on $\mathbb{R}$ with finite second moment, the high-rate
distortion of the optimal non-uniform quantizer with $K$ intervals is
\[
D_K^{\mathrm{non-uniform}}
= \frac{1}{12}\,K^{-2}
  \left(\int_{-\infty}^{\infty} f(x)^{1/3}\,dx\right)^3
  + o(K^{-2}), \qquad K \to \infty.
\]
For Gaussian $f$, the integral is finite over $\mathbb{R}$, giving asymptotic
constant $C_{\mathrm{non-uniform}} = \tfrac{1}{12}
\bigl(\int_{-\infty}^\infty f(x)^{1/3}\,dx\bigr)^3$.

\subsection{Uniform Partitioning Error Bound}

A uniform scheme partitions $[-R, R]$ into $K$ equal intervals of width $\Delta =
2R/K$. The distortion is
\[
D_K^{\mathrm{uniform}}
= \frac{1}{12}\,\Delta^2 \sum_{i=1}^K \Pr(X \in I_i)
  + 2\int_R^\infty (x-R)^2 f(x)\,dx.
\]
Choosing $R = \Theta(\sqrt{\log K})$ makes the tail term negligible, yielding
\[
D_K^{\mathrm{uniform}}
= \frac{R^2}{3K^2} + o(K^{-2}),
\]
with constant $C_{\mathrm{uniform}} = R^2/3$.

\subsection{Comparison}

By Holder's inequality with exponents $p = 3$ and $q = 3/2$,
\[
\left(\int_{-\infty}^{\infty} f(x)^{1/3}\,dx\right)^3
< R^2 \int_{-\infty}^{\infty} f(x)\,dx = R^2,
\]
where the inequality is strict for any non-constant $f$. Therefore
\[
C_{\mathrm{non-uniform}}
= \frac{1}{12}\left(\int_{-\infty}^{\infty} f(x)^{1/3}\,dx\right)^3
< \frac{R^2}{12}
< C_{\mathrm{uniform}},
\]
and hence $D_K^{\mathrm{ASNC}} < D_K^{\mathrm{SpikeZIP-FT}}$ asymptotically. This
completes the proof of Theorem~\ref{thm:asnc_bound} in the main paper.

\begin{remark}
The extension from the nonnegative half-axis to the full real line is handled
naturally here. For the zero-crossing region where $a_i < 0 < b_i$, the
morphology-aware weighting assigns regional priority $R_{\mathrm{region}} = 3.0$
(Appendix~\ref{app:asnc_details}), ensuring denser partitioning in the region of
highest curvature for SiLU and related activations. This is consistent with the
optimal non-uniform allocation predicted by the Bennett integral, which concentrates
intervals where $f(x)^{1/3}$ is largest.
\end{remark}

\subsection{ASNC Error Localization}
\label{app:error_localization}

\begin{corollary}[ASNC error localization]
\label{cor:asnc_localization}
Let $\epsilon_{\mathrm{ASNC}}$ denote the approximation error at a nonlinear module.
The ASNC output is immediately re-encoded into BSE format
via~\eqref{eq:asnc_output}. By Corollary~\ref{cor:no_diffusion}, all subsequent
linear layers operate on this BSE-encoded output exactly, introducing no additional
error beyond the quantization floor. Therefore $\epsilon_{\mathrm{ASNC}}$ does not
accumulate or amplify: the total network error with $M$ nonlinear modules is bounded
by $M \cdot \epsilon_{\mathrm{ASNC}}$, independent of network depth $L$.
\end{corollary}

\begin{proof}
By construction~\eqref{eq:asnc_output}, the ASNC output $\mathbf{S}^{\mathrm{BSE}}$
is a valid BSE spike train. By Corollary~\ref{cor:no_diffusion}, any linear layer
receiving a valid BSE input produces a valid BSE output with no additional error
beyond the quantization floor. The ASNC approximation error is therefore frozen at
the point of re-encoding and cannot be amplified by subsequent exact linear
operations. Summing over $M$ nonlinear modules gives the stated bound.
\end{proof}

\begin{remark}
This contrasts sharply with prior methods, where per-layer error $\epsilon$ compounds
to $O((1+\epsilon)^L)$ after $L$ layers. Under LASER the bound is
$O(M \cdot \epsilon_{\mathrm{ASNC}})$, linear in $M$ and independent of $L$. For a
typical transformer where $M \ll L$, this represents an exponential-to-linear
improvement in error growth with depth.
\end{remark}


%%%%%%%%%%%%%%%%%%%%%%%%%%%%%%%%%%%%%%%%%%%%%%%%%%%%%%%%%%%%%%%%%%%%%%%%%%%
\section{STE Gradient Characterization Under LASER}
\label{app:ste_proof}
%%%%%%%%%%%%%%%%%%%%%%%%%%%%%%%%%%%%%%%%%%%%%%%%%%%%%%%%%%%%%%%%%%%%%%%%%%%

This appendix provides the full proof of Proposition~\ref{prop:ste_main} in the main
paper. We restate the result in the notation of this appendix for completeness.

Standard surrogate gradient methods apply STE on top of lossy representations, where
the forward-backward inconsistency has no computable error bound. Under LASER, the
high-fidelity forward path enables a quantitative characterization. We emphasize that
the result characterizes the gradient of the \emph{$N$-bit quantized ANN} for linear
layers, not the gradient of the full-precision ANN; the two differ by the quantization
gradient, which is a separate and standard concern in quantization-aware training.

\begin{proposition}[STE gradient characterization under LASER, full statement]
\label{prop:ste_correct}
Under Assumptions~\ref{asm:quantization_alignment} and~\ref{asm:input_spacing}:
\begin{enumerate}[label=(\roman*),leftmargin=*]
  \item \textbf{Linear layers.} The STE gradient equals the true gradient of the
        $N$-bit quantized ANN forward pass. No approximation is introduced beyond the
        ANN's quantization gradient.
  \item \textbf{Nonlinear layers.} Under Assumption~\ref{asm:input_spacing}, the STE
        gradient error relative to the true gradient of $\sigma$ at the input is
        bounded by $\epsilon_{\mathrm{ASNC}} / \delta_{\min}$, where
        $\delta_{\min} = \min_i \delta_i > 0$ is the minimum input spacing across
        segments and $\epsilon_{\mathrm{ASNC}} \le D_K^{\mathrm{ASNC}}$.
  \item \textbf{Depth independence.} Under the no-amplification property of
        Corollary~\ref{cor:asnc_localization}, the total gradient error sums rather
        than compounds: it is bounded by $O(M \cdot \epsilon_{\mathrm{ASNC}} /
        \delta_{\min})$, independent of network depth $L$.
\end{enumerate}
\end{proposition}

\begin{proof}
\textbf{(i) Linear layers.}
By Corollary~\ref{cor:no_diffusion}, the BSE forward pass through any linear layer
$f(\mathbf{x}) = \mathbf{W}\mathbf{x} + \mathbf{b}$ is functionally identical to the
$N$-bit quantized ANN forward pass. The gradient of the quantized ANN with respect
to its input is $\mathbf{W}^\top$ (the quantization step itself is piecewise constant
and non-differentiable, but the STE treats it as the identity, consistent with
quantization-aware training practice). The STE for the BSE encoder also substitutes
the identity. Since the BSE encoder is the identity on the integer grid by
Theorem~\ref{thm:bijective}, the composition of the two STE substitutions leaves the
chain rule through the linear layer unaffected, and the STE gradient equals the
gradient of the quantized ANN.

\textbf{(ii) Nonlinear layers.}
At an ASNC module approximating $\sigma$, the forward pass introduces an error
$\epsilon_{\mathrm{ASNC}} = |\hat{\sigma}(x) - \sigma(x)| \le D_K^{\mathrm{ASNC}}$
in the $L^2$ sense. The STE treats the ASNC module as an identity in the backward
pass, passing the upstream gradient $g$ through unchanged. The true backward signal
would be $g \cdot \sigma'(x)$; the STE delivers $g$. The difference is
$g \cdot (1 - \sigma'(x))$.

Now, the ASNC approximation error $\epsilon_{\mathrm{ASNC}}$ bounds the deviation of
$\hat{\sigma}$ from $\sigma$. Within segment $i$, by the mean value theorem applied
to $\hat{\sigma} - \sigma$, the derivative error satisfies
$|\hat{\sigma}'(x) - \sigma'(x)| \le \epsilon_{\mathrm{ASNC}} / \delta_i$
where $\delta_i > 0$ is the minimum input spacing in segment $i$
(Assumption~\ref{asm:input_spacing}). Since the STE gradient corresponds to using
$\hat{\sigma}' = 0$ (identity substitution) rather than $\sigma'(x)$, the gradient
error relative to the true gradient of $\hat{\sigma}$ is bounded by
$|g| \cdot |\hat{\sigma}'(x)| \le |g| \cdot \epsilon_{\mathrm{ASNC}} / \delta_{\min}$.

\textbf{(iii) Depth independence.}
By Corollary~\ref{cor:asnc_localization}, $\epsilon_{\mathrm{ASNC}}$ from one
nonlinear module is not amplified by subsequent linear layers — BSE encodes the
output at the quantization floor, preventing further growth. Each nonlinear module
contributes a bounded gradient error independently; these contributions sum rather
than multiply. The total gradient error is therefore
$O(M \cdot \epsilon_{\mathrm{ASNC}} / \delta_{\min})$, independent of $L$.
\end{proof}

\begin{remark}
Proposition~\ref{prop:ste_correct} provides a quantitative characterization of STE
gradient error under LASER. In prior work, STE is applied on top of lossy
representations where the forward-backward inconsistency has no computable bound.
Under LASER, the inconsistency is equal to the quantized ANN gradient for linear
layers and bounded by $\epsilon_{\mathrm{ASNC}} / \delta_{\min}$ for nonlinear
layers. The result is a characterization, not a proof that STE is exact everywhere:
at nonlinear layers, a genuine bounded error remains, and its magnitude depends on
both $K$ (controlled by the practitioner via ASNC resolution) and $\delta_{\min}$
(controlled by the segment calibration). Increasing $K$ tightens $D_K^{\mathrm{ASNC}}$
but may decrease $\delta_{\min}$, creating a tradeoff that should be tuned in
practice.
\end{remark}


%%%%%%%%%%%%%%%%%%%%%%%%%%%%%%%%%%%%%%%%%%%%%%%%%%%%%%%%%%%%%%%%%%%%%%%%%%%
\section{ASNC Formal Specification}
\label{app:asnc_details}
%%%%%%%%%%%%%%%%%%%%%%%%%%%%%%%%%%%%%%%%%%%%%%%%%%%%%%%%%%%%%%%%%%%%%%%%%%%

\subsection{Segmental Codec and BSE Re-encoding}

Each segmental codec $\mathcal{C}_i(x)$ uses a spiking neural structure as defined
in~\eqref{eq:asnc_codec} of the main paper. The decoded analog value $\hat{y}_i$ for
each active segment is immediately re-encoded by the BSE encoder to produce
$\mathbf{S}^{\mathrm{BSE}}_i(t)$ within the same window $T$; the overall output
follows~\eqref{eq:asnc_output}.

\subsection{Adaptive Splitting Algorithm}

\noindent\textbf{Split criterion.}
\begin{equation}
\frac{L_i}{\bar{L}} \cdot \big(1 + \sigma_i^2\big) \cdot \rho_i
\;>\; \tau_{\mathrm{adaptive}},
\end{equation}
where $L_i$ is the per-segment loss, $\bar{L}$ is the global average loss,
$\sigma_i^2$ is the loss variance within the segment, and $\rho_i$ is the activation
proportion.

\noindent\textbf{Adaptive threshold.}
\begin{equation}
\tau_{\mathrm{adaptive}}
\;=\; \tau_{\mathrm{base}} \cdot \alpha_{\mathrm{stag}} \cdot \alpha_{\mathrm{prog}},
\end{equation}
where $\tau_{\mathrm{base}}$ is a base threshold and the $\alpha$ factors correct for
training stagnation and progress:
\begin{equation}
\alpha_{\mathrm{stag}} \;=\;
\begin{cases}
0.2, & \text{if } N_{\mathrm{stag}}/N_{\mathrm{active}} > 0.6,\\
0.4, & \text{if } N_{\mathrm{stag}}/N_{\mathrm{active}} > 0.3,\\
1.0, & \text{otherwise,}
\end{cases}
\qquad
\alpha_{\mathrm{prog}} \;=\;
\begin{cases}
0.4, & \text{if } (L_0-L_t)/L_0 < 0.005,\\
0.6, & \text{if } (L_0-L_t)/L_0 < 0.02,\\
1.0, & \text{otherwise.}
\end{cases}
\end{equation}

\noindent\textbf{Split point selection.}
Candidate split points $x_s$ are selected by maximizing
\begin{equation}
S(x_s)
= \alpha_1\,E_{\mathrm{sep}}(x_s)
+ \alpha_2\,C_{\mathrm{cons}}(x_s)
+ \alpha_3\,B_{\mathrm{bal}}(x_s)
+ \alpha_4\,Y_{\mathrm{sep}}(x_s),
\end{equation}
where $E_{\mathrm{sep}}$ measures error separation, $C_{\mathrm{cons}}$ reflects
internal consistency, $B_{\mathrm{bal}}$ evaluates data balance, and $Y_{\mathrm{sep}}$
captures target separation.

\subsection{Adaptive Weighting Mechanism}

\noindent\textbf{Morphology-aware weights.}
\begin{equation}
w_i = 0.5\,I_{\mathrm{func}}(i) + 0.35\,I_{\mathrm{error}}(i) + 0.15\,D_i,
\end{equation}
where $I_{\mathrm{func}}(i)$ is the functional importance of segment $i$,
$I_{\mathrm{error}}(i)$ quantifies its error profile, and $D_i$ is its data density.

\noindent\textbf{Regional priority $R_{\mathrm{region}}(i)$.}
\begin{equation}
R_{\mathrm{region}}(i) \;=\;
\begin{cases}
3.0, & a_i < 0 \text{ and } b_i > 0 \quad \text{(zero-crossing region)},\\
2.5, & b_i \le 0 \quad \text{(purely negative region)},\\
2.0, & a_i \ge 0 \quad \text{(purely positive region)},\\
1.0, & \text{otherwise.}
\end{cases}
\end{equation}

\end{document}